\RequirePackage{fix-cm}
\documentclass[11pt,a4paper]{ltjsarticle}
\usepackage[margin=24mm]{geometry}
\usepackage{amsmath,amssymb,amsthm,mathtools,booktabs}
\usepackage[haranoaji]{luatexja-preset}
\usepackage[unicode,hidelinks]{hyperref}
\usepackage{bookmark}
\setlength{\emergencystretch}{3em}
\newtheorem{theorem}{定理}[section]
\newtheorem{proposition}[theorem]{命題}
\newtheorem{lemma}[theorem]{補題}
\newtheorem{corollary}[theorem]{系}
\theoremstyle{definition}\newtheorem{remark}[theorem]{注意}
\renewcommand{\proofname}{証明}
\newcommand{\OO}{\mathcal O}\newcommand{\cA}{\mathcal A}\newcommand{\cB}{\mathcal B}
\DeclareMathOperator{\Hom}{Hom}
\DeclareMathOperator{\Ext}{Ext}
\DeclareMathOperator{\End}{End}
\DeclareMathOperator{\Perf}{Perf}
\DeclareMathOperator{\RHom}{RHom}
\DeclareMathOperator{\Pic}{Pic}
\DeclareMathOperator{\qgr}{qgr}
\DeclareMathOperator{\coh}{coh}
\DeclareMathOperator{\Proj}{Proj}
\DeclareMathOperator{\Spec}{Spec}
\DeclareMathOperator{\Ind}{Ind}
\DeclareMathOperator{\Jac}{Jac}
\newcommand{\Db}{D^b}
\title{非可換del Pezzo曲面のheart：共鳴根の拡大と三次曲面corner\\
\large 符号別Ext・形式変形・有限次元代数の完全判定}
\author{}
\date{2026年11月17日\qquad 公表準備稿 hearts2-v1}
\hypersetup{pdftitle={非可換del Pezzo曲面の共鳴heartと三次曲面corner}}
\begin{document}
\maketitle
\begin{abstract}
本稿は主論文\texttt{dp-hearts.tex}の自然heart比較と全ADE有限商を背景に，そこに詳細を収録しなかった二つの方向を研究する．
第一にRains共鳴における根sheafの符号別Ext，正負の拡大の最大rank，混合$A_2$の対象と表現，最大共鳴点でのADE zigzag自己Extと固定dg圏内のframed形式変形を解析する．
第二に三次曲面型rolled-up quiverの一頂点除去商について，四交比トーラス上の十一条件による有限次元性の必要十分判定，六点配置との同型，一様Hilbert関数と階数162の有限平坦族，厳密な有理数計算証明書を与える．
最大共鳴の全次数でのZariski稠密性，零指数壁の収縮と非可換弱gluingのorbit核，規格化AS heartの自然比較は主論文で証明しているので，本稿では必要な入力だけを再掲する．
基礎的なzigzag形式性やgeometric helixの既知結果と，本研究固有の拡大・corner計算の寄与を区別する．
\end{abstract}
\tableofcontents

\part{Rains共鳴の符号別Ext・混合拡大と形式変形}
\section{共鳴の規約と入力の範囲}
Rainsの既知のeffective-root理論，境界triangleと反標準環の構成を入力とする．
zigzagの形式性はLiu--Wangの既知定理であり，本研究の新定理として扱わない．
全$q$の導来比較は主論文\cite{Hearts}の\texttt{dp-hearts.tex}に証明がある．
主論文ではさらに，全ADE境界外heartのpreprojective模型と，
元の$A$の全有限次元表現を支配する商を扱う．本稿は符号別Ext・混合拡大・形式変形と，三次曲面のcornerの完全計算を主題とする．

基礎体を$k=\mathbb C$とする．Rainsの意味のsmooth elliptic
noncommutative del Pezzo $(\mathcal A_{\mathrm R},E)$に
構造line bundle $e=\mathcal O_X$，Serre関手$\mathbb S$，
$\theta=\mathbb S[-2]$，$\tau=\theta^{-1}$，
境界制限$Li^*:\Db(\mathcal A_{\mathrm R})\to\Perf(E)$
と右随伴$i_*$があり，
\begin{equation}\label{eq:boundarytriangle}
 \theta M\longrightarrow M\longrightarrow i_*Li^*M
 \longrightarrow\theta M[1]
\end{equation}
は自然なtriangleである\cite[\S\S4--5]{RainsFiltered}．
$e$はline bundleで$\theta,\tau$はRains heartにexactに作用する．
さらに$R=\bigoplus_{n\ge0}\Hom(e,\tau^ne)$は
次数1の中心元$g$を持ち，$R/(g)$は楕円曲線上のtwisted
homogeneous coordinate ringとなる\cite[Proposition 5.1]{RainsFiltered}．
以下，$\qgr R$の対象と有限生成次数付き加群を扱う場合には，
当該$R$の右noetherian性と必要な$\chi$条件を仮定または確認した範囲に限る．

古典的avatarを滑らかなgenuine del Pezzo曲面$S$，
その滑らかな反標準曲線を$E$とし，
$K_0(\mathcal A_{\mathrm R})\cong K_0(S)$を固定する．
根$\alpha\in K_S^\perp$は$\alpha^2=-2$を満たす．
$M$のrankは零，$c_1(M)=\alpha$，
$m=\chi(e,M)$とする．
gluingパラメータ$q\in\Pic^0(E)$の向きを
\begin{equation}\label{eq:resonance}
 Li^*M=0\quad\Longrightarrow\quad
 \mathcal O_E(\alpha)\otimes q^{m}\simeq\mathcal O_E
\end{equation}
と規格化する．これは$K$群上のrestriction determinant
の式である．Rainsの$q$を逆元とする規約では指数の符号も反転する．
式\eqref{eq:resonance}だけではsheafの存在を保証しない．
以下では，Rainsの既約根sheaf存在定理
\cite[\S9, \S11, pp.~139--141]{RainsBirational}と
\cite[\S4, pp.~33--35]{RainsFiltered}が適用され，
$Li^*M=0$である\emph{既約な純一次元sheaf}が実在する共鳴を扱う．
特に，固定したblowdownのsimple rootの共鳴を使えばよい．
非torsion仮定を使う単純性・有限長性は\S\ref{sec:mixed-a2-category}で別に証明する．

\begin{proposition}[共鳴による反標準環のHilbert関数の跳びはない]
\label{prop:hilbert-no-jump}
前節のRainsの環の仮定の下，$d=(-K_S)^2>0$とすると
\[
 \dim_k R_n=1+\frac d2 n(n+1)\qquad(n\ge0)
\]
であり，$q$の共鳴の有無に依存しない．
特に非零指数の共鳴で生じる有限次元アフィン表現は，
$R$のHilbert関数の跳びと同一視してはならない．
\end{proposition}
\begin{proof}
$g$は正則中心元であり
$R/(g)=B(E,L,\sigma)$で$\deg L=d$だから
\[
0\longrightarrow R_{n-1}\xrightarrow{\cdot g}R_n
\longrightarrow B(E,L,\sigma)_n\longrightarrow0
\]
が成立する．$n\ge1$では楕円曲線上の正次数線束の
切断数は$nd$である．$R_0=k$から帰納法で主張が従う．
このgraded ringの恒等式はRains
\cite[Proposition 5.1]{RainsFiltered}の仮定内での結果である．
\end{proof}

\section{根sheafのEuler指数とcohomological degree}
\label{sec:sign-trichotomy}
Rainsの$\mathrm{ad}$は，一般に同じ非可換曲面の自己双対ではなく，
反対側の双対曲面$X^{\mathrm{ad}}$への反変導来同値である．
構造line bundle $e$はその双対$e^{\mathrm{ad}}=\mathrm{ad}(e)$へ移る．
既知の双対性
\[
 \mathbf R\mathrm{ad}_X:
 D^b(\mathcal A_{\mathrm R})^{\mathrm{op}}
 \xrightarrow{\sim}D^b(\mathcal A_{\mathrm R}^{\mathrm{ad}})
\]
とその二重双対を使う\cite[\S5, Corollary 5.9]{RainsBirational}．

\begin{lemma}[純一次元の双対と指数反転]
\label{lem:root-dual}
$M$を境界に制限して零となる既約な純一次元sheafとする．
$N=R^1\mathrm{ad}_X(M)$は双対曲面上の既約な純一次元sheafで，
\[
 \mathbf R\mathrm{ad}_X M\simeq N[-1],\quad
 R^1\mathrm{ad}_{X^{\mathrm{ad}}}N\simeq M,\quad
 c_1(N)=c_1(M),\quad \chi(e^{\mathrm{ad}},N)=-\chi(e,M).
\]
$M$上で自然な$\theta M\to M$が同型なら，
$N$上でも自然な$\theta^{\mathrm{ad}}N\to N$が同型である．
とくに，根sheafは双対でも根sheafになる．
\end{lemma}
\begin{proof}
Rains\cite[Corollaries 7.26, 7.30, 7.31]{RainsBirational}
によれば純一次元sheafでは$\mathrm{ad}$の導来cohomologyが
次数1だけにあり，その二重双対は元のsheafを復元する．
これにより$N$も純一次元である．
既約性はgenericな一次元support上の単純加群性に関する性質であり，
genericでの双対によって保たれる．
自然な境界triangleの双対は，双対側の境界triangleに一致するため，
自然変換が$M$上で可逆なら$N$上でも可逆である．
K群では$[\mathbf R\mathrm{ad}M]=-[N]$である．
Rains\cite[Proposition 7.16]{RainsBirational}の公式
$c_1(\mathbf R\mathrm{ad}M)=-c_1(M)$，
$\chi(\mathbf R\mathrm{ad}M)=\chi(M)$を使えば
$c_1(N)=c_1(M)$，$\chi(N)=-\chi(M)$となる．
ここで$\chi(N)=\chi(e^{\mathrm{ad}},N)$である．
\end{proof}

\begin{lemma}[零以下のEuler指数ではglobal sectionがない]
\label{lem:no-section}
$M$を根$\alpha^2=-2$を第一Chern類とする既約な純一次元sheafとする．
$m=\chi(e,M)\le0$なら$\Hom(e,M)=0$．
\end{lemma}
\begin{proof}
非零の$f:e\to M$があるとし，その像を$I$とする．
$M$の既約性は「任意の非零部分sheafが同じ第一Chern類を持つ」
ことであって，$f$が全射であることを意味しない
\cite[\S11, pp.~138--140]{RainsBirational}．
$I\ne0$だから$c_1(I)=\alpha$であり，商$T=M/I$は
零または零次元である．
$I$は$e$の商であるため，
Rains\cite[Corollary 10.16]{RainsBirational}の不等式
\[
 \chi(e,I)\ge-\tfrac12\alpha(\alpha+K_X)=1
\]
が適用できる．
一方，零次元sheafのEuler標数は非負
\cite[Proposition 7.6]{RainsBirational}なので
$\chi(e,M)=\chi(e,I)+\chi(e,T)\ge1$となり，
$m\le0$に矛盾する．
\end{proof}

\begin{theorem}[resonance sheafの符号別Extの完全判定]
\label{thm:sign}
$M$は$\alpha^2=-2$，$(-K_X)\alpha=0$を満たす根の
既約な純一次元sheafで，$Li^*M=0$とする．
$m=\chi(e,M)\in\mathbb Z$について全ての$n\in\mathbb Z$で
\[
\boxed{
\begin{aligned}
\dim\Hom(e,\tau^nM)&=\max(m,0),\\
\dim\Ext^1(e,\tau^nM)&=\max(-m,0),\\
\Ext^2(e,\tau^nM)&=0.
\end{aligned}}
\]
したがって自然な導来切断を$M$に適用すると，
\begin{equation}\label{eq:trichotomy}
 \Phi(M)=
 \begin{cases}
   \pi N_m,&m>0,\\
   0,&m=0,\\
   (\pi N_{-m})[-1],&m<0,
 \end{cases}
\end{equation}
ただし$m\ne0$のとき$\pi N_{|m|}\in\qgr R$は
境界$g=0$の外にある非零対象であり，
$R/(g-1)$の$|m|$次元右加群に対応する．
$m=0$のとき，$M$は非零であるが，全Serre orbitに
導来的に右直交し，指定した反標準導来切断は$M$を零に送る．
\end{theorem}
\begin{proof}
$Li^*M=0$なら境界triangleから自然な
$\theta M\to M$が同型であり，$\tau^nM\simeq M$．
line bundle $e$と一次元sheafに対する
$\Ext^2(e,M)=0$は
Rains\cite[Corollary 7.23]{RainsBirational}による．

$m\le0$では補題\ref{lem:no-section}により
$\Hom(e,M)=0$，Euler標数より$\dim\Ext^1(e,M)=-m$．
特に$m=0$では全$\Ext$が零となる．

$m>0$とする．双対$N=R^1\mathrm{ad}M$は
補題\ref{lem:root-dual}により双対曲面上の
負指数$-m<0$の既約根sheafで，境界にも制限して零である．
従って$\Hom_{X^{\mathrm{ad}}}(e^{\mathrm{ad}},N)=0$，
$\dim\Ext^1_{X^{\mathrm{ad}}}(e^{\mathrm{ad}},N)=m$．
Rainsの反変導来双対とSerre双対性を合成すると
\begin{equation}\label{eq:dual-ext}
\begin{aligned}
\Hom_X(e,M)
 &\cong\Ext^1_{X^{\mathrm{ad}}}(e^{\mathrm{ad}},N)^*,\\
\Ext^1_X(e,M)
 &\cong\Hom_{X^{\mathrm{ad}}}(e^{\mathrm{ad}},N)^*.
\end{aligned}
\end{equation}
実際$\mathbf R\mathrm{ad}_XM=N[-1]$なので
\[
 \Ext^{j}_{X^{\mathrm{ad}}}(e^{\mathrm{ad}},N)
 \cong\Ext_X^{j+1}(M,e)
 \cong\Ext_X^{1-j}(e,\theta M)^*
\quad (j=0,1).
\]
これに$\theta M\simeq M$を代入すれば
$\dim\Hom_X(e,M)=m$，$\Ext_X^1(e,M)=0$を得る．

最後に$g$は$\theta\to\mathrm{id}$から生じる中心正則元で，
$M$上で同型として作用するので，
各次数の$\Hom$または$\Ext^1$の$g$による移行は同型．
非零の次数付き$R$加群$N_{|m|}$は
$k[g]$上有限生成で，各次数の次元が$|m|$，
局所化の次数0部分は$|m|$次元の$R/(g-1)$右加群になる．
非零cohomologyが一つの次数に集中しているので
\eqref{eq:trichotomy}が従う．
\end{proof}

\begin{proposition}[root sheafは2-spherical]
\label{prop:spherical}
定理\ref{thm:sign}の$M$は任意の$m$について
\[
\Ext^j(M,M)=
\begin{cases}k&j=0,2,\\0&j\ne0,2.\end{cases}
\]
特に$\theta M\simeq M$の下で$2$-spherical objectである．
\end{proposition}
\begin{proof}
$M$の非零自己準同型の像は既約性により
第一Chern類$\alpha$を持つ．
核が非零なら零次元であり，$M$の純一次元性に反する．
よって自己準同型は単射であり，その余核は零次元．
自己射なので余核の数値類は零，Rains
\cite[Proposition 7.6]{RainsBirational}から余核も零．
従って$\End(M)$は有限次元の斜体で，$k=\mathbb C$だから
$\End(M)=k$．Serre双対により
$\Ext^2(M,M)=\End(M)^*=k$．
Rains\cite[Corollary 7.13]{RainsBirational}のEuler形式で
$\chi(M,M)=-\alpha^2=2$となるので$\Ext^1(M,M)=0$．
\end{proof}

\begin{corollary}[pointedなheart非一致と導来切断の限界]
\label{cor:pointed-sign}
$m<0$の既約根sheafがあると，Rains heartの$(e,\tau)$は
AZ-ampleでない．標準pair$(\pi R,(1))$がAZ-ampleなら
$(\mathcal A_{\mathrm R},e,\tau)$と
$(\qgr R,\pi R,(1))$はpointed abel圏として同値でない．
さらに指定した自然な$\Phi$は$m<0$で$t$-exactでなく，
$m=0$で保守的でない．
$m>0$の個々の根sheafの像はheartの次数0にあるが，
そのことだけでは$\Phi$の全体の$t$-exactnessは従わない．
\end{corollary}
\begin{proof}
$m<0$では非零$M$に対して全ての
$\Hom(e,\tau^nM)$が零なので，AZの生成公理が失敗する．
pointedなabel同値はample pairを保存する．
後二者は\eqref{eq:trichotomy}から従う．
\end{proof}
\begin{remark}[先行する証明の訂正と範囲]
「既約root sheafへの非零な構造sheafからの射は全射」
という推論は正しくない．必要なのは
像が同じ第一Chern類を持ち，商が零次元であることだけである．
補題\ref{lem:no-section}がその正しい証明を与える．
$m=0$の対象が実在することは共鳴の数値条件からは従わない．
genuine del Pezzo avatarでは後述の命題により$m=0$の根共鳴はない．
また指定した$\Phi$の非忠実性は，
別の抽象的導来同値の不存在を意味しない．
\end{remark}

\section{境界外のsheafと拡大によるcohomologyの相殺}
\label{sec:offboundary-extensions}
本節は既約根sheafに限らない．既知入力は
Rainsの境界triangle \eqref{eq:boundarytriangle}，
$1$次元以下のsheafへの$\Ext^2$消滅
\cite[Corollary 7.23]{RainsBirational}，および
中心正則元$g$を持つ反標準環の仮定である．
$\mathcal Z$をRains heartの中で
\[
 \mathcal Z=\{M\in\mathcal A_{\mathrm R}:
    \dim M\le1,\ Li^*M\simeq0\}
\]
と置く．$R$は右noetherianであり，必要な$\chi$条件により
$\Phi=\pi\mathbf R\Gamma_*$を該当対象に定義できると仮定する．
$A=R/(g-1)\cong (R[g^{-1}])_0$と書く．
ここで$\pi$は有限次元次数付き加群によるSerre商である．

\begin{theorem}[境界外のsheafに対する二項振幅と長完全列]
\label{thm:offboundary-amplitude}
$\mathcal Z$は拡大に閉じたexact categoryである．
任意の$M\in\mathcal Z$について次が成立する：
\begin{enumerate}
\item 自然な$\theta M\to M$は同型である．従って
$\tau^nM\simeq M$で，全$n\in\mathbb Z$について
$\Ext^j(e,\tau^n M)=0$ $(j\notin\{0,1\})$．
\item
\[
 N^j(M)=\bigoplus_{n\ge0}\Ext^j(e,\tau^nM)
 \qquad(j=0,1)
\]
は次数付き右$R$加群で，$g$による乗法は全次数で同型である．
特に$k[g]$加群として有限生成freeで，そのrankは
$h^j(M)=\dim_k\Ext^j(e,M)$．
\[
 V^j(M)=(N^j(M)[g^{-1}])_0\in\operatorname{mod}_{fd}A,
 \qquad \dim_kV^j(M)=h^j(M)
\]
であり，$H^j(\Phi M)\simeq\pi N^j(M)$．
従って自然な導来切断のcohomologyは次数$0,1$だけで，
どちらも反標準境界$g=0$の外にある．
\item 任意の短完全列$0\to M_1\to M_2\to M_3\to0$
（$M_i\in\mathcal Z$）から，有限次元右$A$加群の
自然な完全列
\begin{equation}\label{eq:offboundary-six}
0\to V^0(M_1)\to V^0(M_2)\to V^0(M_3)
\xrightarrow{\delta} V^1(M_1)
\to V^1(M_2)\to V^1(M_3)\to0
\end{equation}
を得る．この$\delta$はextension classのconnecting morphismで，
$A$-線型である．
\end{enumerate}
\end{theorem}
\begin{proof}
Rainsのdimension filtrationは短完全列に関し安定である．
また$Li^*$はexact triangulated functorなので
両端の制限が零なら中間項の制限も零となる．
従って$\mathcal Z$は拡大に閉じる．

境界triangle \eqref{eq:boundarytriangle}と$Li^*M=0$から
自然な$\theta M\to M$が同型となる．
$\theta$，$\tau$はheart上でexactなので，全twistも
$1$次元以下のsheafとなる．
負のExtは二つのheart対象間で消え，Rainsの
\cite[Corollary 7.23]{RainsBirational}で$\Ext^2$が消える．
この曲面のcohomological dimensionは高々$2$だから，
(1)が従う．

境界triangleの$\theta\to\mathrm{id}$から得られる
次数$1$の中心元$g$は$M$上で同型であり，
$\Ext^j(e,\tau^nM)\to\Ext^j(e,\tau^{n+1}M)$に
同型を誘導する．有限次元性により
$N^j(M)\simeq k[g]\otimes_k\Ext^j(e,M)$
が$k[g]$加群として成り立つ（$R$作用はこの分解で
単純になるとは主張しない）．
ゆえに$\pi N^j$は非零なら境界外の対象で，
dehomogenizationは$A$の有限次元加群$V^j$を与える．
$\Phi M$のcohomologyの同定は$\pi$のexactnessと
$\mathbf R\Gamma_*$の定義による．

短完全列の$\tau^n$-twistに$\RHom(e,-)$を適用すると
cohomologyは次数$0,1$だけなので六項完全列になる．
これを全$n\ge0$で直和し，中心元$g$で局所化した後
次数$0$を取る．局所化と次数$0$の抽出はこの
graded-moduleのクラスでexactであり，
自然性によりconnecting morphismは$A$-線型となる．
\end{proof}

\begin{corollary}[拡大に閉じた正・負の部分圏]
\label{cor:offboundary-sign-exact}
\[
 \mathcal Z_+=\{M\in\mathcal Z:\Ext^1(e,M)=0\},\qquad
 \mathcal Z_-=\{M\in\mathcal Z:\Hom(e,M)=0\}
\]
はともに拡大に閉じる．
$\Phi(\mathcal Z_+)$はanticanonical heartの次数$0$に，
$\Phi(\mathcal Z_-)$はその$[-1]$ shiftに含まれる．
$V^0:\mathcal Z_+\to\operatorname{mod}_{fd}A$，
$V^1:\mathcal Z_-\to\operatorname{mod}_{fd}A$
はそれぞれ短完全列を保つ．
特に正指数（負指数）の既約根sheafの任意の有限回の
拡大は$\mathcal Z_+$（$\mathcal Z_-$）に入る．
\end{corollary}
\begin{proof}
$\Ext^1$，$\Hom$の各々の消滅は
\eqref{eq:offboundary-six}で拡大に関して閉じる．
定理\ref{thm:sign}により各既約根sheafは指定部分圏に属する．
長完全列の端の項が零なら残る三項は短完全列となる．
\end{proof}

\begin{proposition}[正負の混合拡大と相殺]
\label{prop:mixed-sign-cancellation}
$m_+>0$，$m_-<0$を持つ境界外既約根sheaf
$M_+\in\mathcal Z_+$，$M_-\in\mathcal Z_-$が同じ
Rains heartに実在すると仮定する．
$0\to M_-\to M\to M_+\to0$という拡大を取る．
この拡大が誘導する$A$-線型写像
\[
 \delta_M:V^0(M_+)\longrightarrow V^1(M_-)
\]
のrankを$r$とすると，
\begin{equation}\label{eq:mixed-ranks}
 \dim_kV^0(M)=m_+-r,\qquad
 \dim_kV^1(M)=-m_--r .
\end{equation}
特に$m_+=-m_->0$の場合，$\delta_M$が同型なら
$\Phi(M)=0$である．
従って$\Phi$が$\mathcal Z$上で保守的なら，
非零の混合拡大に対するそのような同型
$\delta_M$は存在しない．
同じ結論は$\Phi$が全導来圏の同値であるなら必ず成立する．
\end{proposition}
\begin{proof}
\eqref{eq:offboundary-six}に
$V^0(M_-)=V^1(M_+)=0$を代入すると
\[
0\to V^0(M)\to V^0(M_+)
\xrightarrow{\delta_M}V^1(M_-)\to V^1(M)\to0
\]
となり，\eqref{eq:mixed-ranks}が従う．
$\delta_M$が同型なら$N^0(M),N^1(M)$は
次数$0$の局所化が零であるが，$g$は各次数において
同型なので両方零であり，$\Phi(M)=0$である．
$M$は$M_-$を非零部分対象とするため非零である．
\end{proof}

\begin{remark}[Euler標数だけでは検出できない]
一般の$M\in\mathcal Z$に対して
\[
 \chi(e,M)=\dim V^0(M)-\dim V^1(M)
\]
であるが，$\chi(e,M)=0$から$\Phi(M)=0$は従わない．
正負の拡大類によるconnecting morphism
$\delta_M$のrankを指定して初めて
反標準切断のcohomologyが決まる．
$\mathcal Z_+,\mathcal Z_-$が$\mathcal Z$全体の
torsion pairになるとは証明していない．
\end{remark}

\begin{proposition}[根系からの共鳴交叉制約]
\label{prop:root-intersection}
Rains heartに境界に制限して零となる
異なる二つの既約純一次元根sheaf$M_\alpha,M_\beta$が
存在し，第一Chern類を$\alpha\ne\beta$とする．
すると
\[
 \Hom(M_\alpha,M_\beta)=\Ext^2(M_\alpha,M_\beta)=0,
 \qquad
 \dim_k\Ext^1(M_\alpha,M_\beta)=\alpha\cdot\beta.
\]
従って$\alpha\cdot\beta<0$ならそのような二対象は
同じRains heartには同時に存在できず，
$\alpha\cdot\beta=0$なら双方の向きの$\Ext^1$が零である．
$\alpha\cdot\beta=1$なら$\Ext^1$は一次元である．
\end{proposition}
\begin{proof}
非零な$M_\alpha\to M_\beta$の像$I$は，
$M_\beta$の純一次元性から一次元である．
$M_\alpha$の既約性より，核が一次元なら
その第一Chern類は$\alpha$となり，
像は零次元となって矛盾する．
核が零次元なら$M_\alpha$の純一次元性から零である．
従って写像は単射で，像の第一Chern類は$\alpha$である．
他方，$M_\beta$の既約性から非零部分sheafの第一Chern類は
$\beta$である．ゆえに$\alpha=\beta$となり矛盾する．

$\theta M_\alpha\simeq M_\alpha$であるから，
Serre双対により
$\Ext^2(M_\alpha,M_\beta)
 \simeq\Hom(M_\beta,\theta M_\alpha)^*=0$．
Rains\cite[Corollary 7.13, (7.4)]{RainsBirational}の
rank-zero Euler形式は
$\chi(M_\alpha,M_\beta)=-\alpha\cdot\beta$であり，
残る$\Ext^1$の次元は$\alpha\cdot\beta$となる．
\end{proof}

\begin{corollary}[導来完全忠実性に対する有限次元表現テスト]
\label{cor:affine-ext-test}
定理\ref{thm:offboundary-amplitude}の自然な導来切断$\Phi$が
二つの境界外既約根sheaf$M_+,M_-$（指数$m_+>0$，
$m_-<0$）の間の$\Ext^1$で完全忠実ならば，
対応する有限次元$A$加群$V_+$，$V_-$について
\[
\Ext^1_X(M_+,M_-)
 \xrightarrow{\;\sim\;}\Hom_A(V_+,V_-)
\]
が成り立つ．$\alpha_+=c_1(M_+)\ne c_1(M_-)=\alpha_-$なら
\[
\boxed{\dim_k\Hom_A(V_+,V_-)=\alpha_+\cdot\alpha_- .}
\]
従って右辺の次元がこれと異なるなら，
$\Phi$はその二対象の間で完全忠実ではない．
\end{corollary}
\begin{proof}
$\Phi(M_+)\simeq V_+$，$\Phi(M_-)\simeq V_-[-1]$
である（$V_\pm$は境界外の有限支持対象として見る）．
したがって
\[
\Hom_{\Db(\qgr R)}(\Phi M_+,\Phi M_-[1])
 =\Hom_{\qgr R}(V_+,V_-)=\Hom_A(V_+,V_-).
\]
最後の同型は，$g$が各$V_\pm$上で可逆であり，
その次数付き表示が$k[g]$に関して有限自由であることから，
局所化がこの二対象間の$\Hom$を保存するためである．
$\Ext^1_X(M_+,M_-)$は左辺へ$\Phi$が誘導する空間であり，
完全忠実性を仮定すると同型となる．
命題\ref{prop:root-intersection}を使えば次元式を得る．
この写像は命題\ref{prop:mixed-sign-cancellation}の
$\delta_M$のextension classによる線型写像そのものである．
\end{proof}

\begin{proposition}[根の和による反対指数の同時共鳴排除]
\label{prop:root-sum-exclusion}
古典的avatar$S$がgenuine del Pezzoで，
滑らかな反標準楕円曲線$E$を持つとする．
二つの根$\alpha,\beta$に
\[
 \mathcal O_E(\alpha)\otimes q^{m_\alpha}\simeq\mathcal O_E,
 \qquad
 \mathcal O_E(\beta)\otimes q^{m_\beta}\simeq\mathcal O_E
\]
という共鳴があるとする．もし$\alpha\cdot\beta=1$なら，
$m_\alpha+m_\beta=0$は不可能である．
特に，隣接する$A_2$部分根系の根に対する
指数$+r$と$-r$の共鳴は同時には起こらない．
\end{proposition}
\begin{proof}
$\gamma=\alpha+\beta$は
$\gamma^2=-2-2+2=-2$，$H\cdot\gamma=0$を
満たすから再び根である．
$m_\alpha+m_\beta=0$なら二つの共鳴式の積から
$\mathcal O_E(\gamma)\simeq\mathcal O_E$となる．
しかしgenuine del Pezzoについての
命題\ref{prop:zero-index}は任意の根の
$E$への制限が非自明であることを示している．
従って矛盾する．
\end{proof}

\begin{remark}[完全な相殺を探すときの制限]
異なる根sheafの間の非自明な拡大は，
上の命題から$\alpha\cdot\beta>0$を必要とする．
特に最も単純な$\alpha\cdot\beta=1$の場合，
指数が互いに反対で次元も等しいような混合拡大は，
genuine classical avatarではそもそも同時共鳴できない．
これは$\delta_M$の同型による導来的消失が
genuineな族に現れないことの\emph{完全な証明ではない}．
別のroot対や非既約対象を排除できていないためである．
\end{remark}

\begin{remark}[古典的な$(-2)$曲線上の完全な相殺]
\label{ex:euler-cancellation}
滑らかな弱del Pezzo $S$に$(-2)$曲線
$i:D\simeq\mathbb P^1\hookrightarrow S$があり，
反標準曲線$E$は$D$と交わらないとする．
$H=-K_S$の$D$への制限は自明である．
$\mathbb P^1$のEuler完全列
\[
0\to\mathcal O_D(-2)\xrightarrow{(-y,x)}
 \mathcal O_D(-1)^{\oplus2}
 \xrightarrow{(x,y)}\mathcal O_D\to0
\]
を$i_*$で送る．左端を$M_-$，右端を$M_+$
とすると$\chi(M_-)=-1$，$\chi(M_+)=1$である．
中間項$M=i_*\mathcal O_D(-1)^{\oplus2}$は非零だが
\[
\RHom_S(\mathcal O_S(nH),M)
 \simeq\mathbf R\Gamma(D,\mathcal O_D(-1)^{\oplus2})=0
 \quad(n\in\mathbb Z).
\]
したがって$\delta_M:k\to k$は同型で，
反標準導来切断は$M$を零に送る．
これは前の根壁の導来情報消失を
「正負指数の拡大の完全な相殺」として説明する
古典的な模型である．\emph{この例はweak avatarの例であり，
genuine非可換曲面で同じ拡大が存在するという主張ではない．}
\end{remark}


\section{正負の共鳴がそれぞれ稠密}
\begin{theorem}[正負各指数のZariski稠密性]
\label{thm:dense-both}
滑らかな平面三次曲線$E\subset\mathbb P^2$，
無限位数の$q\in\Pic^0(E)$を固定する．
$p_1,p_2,p_4,p_5,p_6\in E$を十分一般に固定し，
$L=\mathcal O_E(1)$と置く．各$n\ge1$と$\varepsilon\in\{+1,-1\}$に対し
\[
 \mathcal O_E(p_3^\varepsilon(n))\simeq
 L\otimes\mathcal O_E(-p_1-p_2)\otimes q^{-\varepsilon n}
\]
を満たす唯一の$p_3^\varepsilon(n)\in E$を取る．
各符号について有限個の例外を除いて六点は一般位置であり，
そのblow-upはgenuineで滑らかなcubic surfaceである．
単純根$\alpha=h-e_1-e_2-e_3$に対し
\[
 \mathcal O_E(\alpha)=q^{\varepsilon n},\qquad
 m=-\varepsilon n.
\]
Rainsの単純根sheaf存在定理の範囲では，
$\varepsilon=+1$の負指数sheafは
反標準切断の$[-1]$に移り，$\varepsilon=-1$の正指数sheafは
次数0にある．それぞれの符号の共鳴点集合が
Zariski denseである．
\end{theorem}
\begin{proof}
$\Pic^1(E)\simeq E$だから$p_3^\varepsilon(n)$は一意に存在する．
$q$が無限位数ゆえ各符号について$p_3^\varepsilon(n)$は相異なる．
五つの固定点に対し，追加の一点が三点共線性や六点
conic上条件を満たすのは有限個の点だけである．
従って一般位置を失う$n$は各符号について有限個であり，
$\{p_3^\varepsilon(n):n\ge1\}$はいずれも
$E$のZariski稠密部分集合である．
根の制限は$L(-p_1-p_2-p_3^\varepsilon(n))=q^{\varepsilon n}$．
式\eqref{eq:resonance}に$m=-\varepsilon n$を代入し，
定理\ref{thm:sign}を適用すれば結論を得る．
\end{proof}

\begin{corollary}[正負共鳴が全パラメータ空間で稠密]
\label{cor:dense-total}
滑らかな平面三次$E$，非torsionの$q\in\Pic^0(E)$，
および他の五吹き上げ点をパラメータとして動かす
六点配置の局所的な族において，定理\ref{thm:dense-both}の正指数および負指数の
共鳴点の合併はそれぞれZariski稠密である．
従ってこのチャートのいかなる非空Zariski開部分にも，
$\mathcal A_{\mathrm R}$で$(e,\tau)$がAZ-ampleでない負指数点が存在する．
\end{corollary}
\begin{proof}
各基底パラメータに対する$p_3$のファイバーは楕円曲線$E$であり，
無限位数$q$を持つ基底点はZariski稠密である．
それらの各ファイバーで定理の正負各指数列はZariski稠密なので，
全空間の任意の非空開集合は適当なファイバーでこの列に交わる．
基底全体のirreducibilityと$p_3$を動かす開写像性を用いる．
\end{proof}
\begin{remark}
ここでいう稠密性は可算個の閉共鳴軌跡の\emph{合併}の
Zariski稠密性であり，個々の共鳴軌跡が開集合であるとか，
全共鳴点で同じheart変化が起こるという主張ではない．
\end{remark}

\begin{remark}[他の次数への同じ稠密性]
$\mathbb P^2$の$r=3,\ldots,8$点blow-up（del Pezzo次数$d=9-r$）
で，滑らかな平面三次$E$上の他の点を一般に固定し，
$p_3(n)$を同じ式で動かせばよい．
genuine del Pezzo条件を破る$p_3$の集合は，固定した他の点が
十分一般なら$E$の有限集合である．
根$\alpha=h-e_1-e_2-e_3$はどの次数でも
$\alpha^2=-2$，$\alpha\cdot K=0$を満たすため，
Rainsの単純根sheaf存在結果の適用範囲で正負各指数の稠密性は
次数$1,\ldots,6$へ同様に及ぶ．これは全moduli成分への
稠密性を主張するものではない．
\end{remark}

\section{実在する混合$A_2$と二段拡大の非消失}\label{sec:dp-hearts2-mixed-a2}
本節の$E_6$最大共鳴の整数表示とZariski稠密性は主論文\cite{Hearts}を入力とする．
数値的共鳴だけから任意の根sheafの存在を推測せず，指定した単純根の実在を確認する．

\begin{proposition}[実在する正負の単純根sheaf]
\label{prop:actual-mixed-a2}
\[
 (N_1,\ldots,N_6)=(1,3,2,5,7,11),\qquad p_i=N_iq
\]
とする．六単純根のEuler指数は
\[
 (m_0,m_1,m_2,m_3,m_4,m_5)=(6,2,-1,3,2,4).
\]
特に
\[
 \alpha=e_1-e_2,\qquad\beta=e_2-e_3,\qquad
 m(\alpha)=2,\quad m(\beta)=-1,\quad
 \alpha\cdot\beta=1.
\]
Rains\cite[\S4, (4.14)--(4.15)]{RainsFiltered}の
simple-root effectivity判定と
\cite[\S9,\S11]{RainsBirational}の既約sheafの存在を
このsmooth elliptic chartに適用すると，同じRains heartに
境界非交叉の既約root sheaf
$M_\alpha,M_\beta$が存在する．それらについて
\[
 \dim\Ext^1(M_\alpha,M_\beta)=1,\quad
 \dim V^0(M_\alpha)=2,\quad
 \dim V^1(M_\beta)=1.
\]
またこの性質を保ち$N_1<N_3<N_2$とする全根共鳴点の集合も
同じ六点配置chartでZariski稠密である．
\end{proposition}
\begin{proof}
単純根指数は$h|_E=3o$と点の差から計算される．
全ての単純根の制限は$\langle q\rangle$にあるので，
Rainsのeffectivity判定によりこれらは有効である．
$E$は滑らかで$\alpha\cdot E=0$なので，
有効な単純根はformal $(-2)$-curveとして
境界と交わらない既約純一次元sheafで実現される．
無限位数の$q$のため$M_\alpha,M_\beta$のEuler指数は
上記の値に一意に固定される．
前節の命題\ref{prop:root-intersection}および
定理\ref{thm:sign}からExtと$V^j$の次元が従う．
最後は主論文\cite{Hearts}の全次数の最大共鳴と整数順序の稠密性による．
\end{proof}

\begin{proposition}[二段混合拡大全体での非消失]
\label{prop:a2-directsum}
命題\ref{prop:actual-mixed-a2}の二sheafを固定する．
$a\ge1,b\ge0$とし，任意の拡大
\[
0\longrightarrow M_\beta^{\oplus b}\longrightarrow M
 \longrightarrow M_\alpha^{\oplus a}\longrightarrow0
\]
を取る．$\Ext^1(M_\alpha,M_\beta)\simeq k$の生成元が
定める$A=R/(g-1)$加群写像を
$d:V_\alpha\to V_\beta$とする．
拡大類に対応する行列$T\in\operatorname{Mat}_{b\times a}(k)$の
connecting morphismは$d\otimes T$であり，
$r=\operatorname{rank}(d)\operatorname{rank}(T)$と置けば
\[
\boxed{
\dim V^0(M)=2a-r,\quad
\dim V^1(M)=b-r,\quad
r\le\min(a,b).}
\]
特に$a\ge1$なら$\dim V^0(M)\ge a>0$なので
自然な反標準導来切断はこの二段拡大を消さない．
$b=2a$では$\chi(e,M)=0$だが，依然として
$\Phi(M)\ne0$である．
\end{proposition}
\begin{proof}
直和とYoneda積の自然性により
\[
\Ext^1(M_\alpha^{\oplus a},M_\beta^{\oplus b})
 \cong\operatorname{Mat}_{b\times a}(k),
\]
対応するconnecting morphismは$k$-線型空間上
Kronecker積$d\otimes T$である．そのrankは積となり，
$d$は二次元から一次元への射だから$\operatorname{rank}(d)\le1$．
前節の六項完全列
\eqref{eq:offboundary-six}に代入して次元公式を得る．
\end{proof}

\begin{corollary}[全$q$導来同値に対する明示的な検査条件]
\label{cor:a2-ff}
同じ例で自然な導来切断$\Phi$が
$M_\alpha,M_\beta$の間の$\Ext^1$で完全忠実なら
\[
 \dim\Hom_A(V_\alpha,V_\beta)=1
\]
である．従って二次元表現$V_\alpha$は
一次元表現$V_\beta$を商に持ち，単純$A$加群ではない．
そのとき非分裂拡大のconnecting morphismはrank$1$で，
$(\dim V^0(M),\dim V^1(M))=(1,0)$となる．
\end{corollary}
\begin{proof}
前節の系\ref{cor:affine-ext-test}による．
$\dim V_\beta=1$だから非零$A$準同型は全射である．
$\Ext^1$の生成元は完全忠実性の下で非零の
$\Hom_A$へ写るから，connecting morphismのrankは$1$．
\end{proof}

\begin{remark}[有理数上の検算例]
\label{rem:a2-rational-witness}
\[
 E:\ y^2z=x^3-xz^2+z^3,\quad o=(0:1:0),\quad
 q=(0:1:1)
\]
とする．楕円曲線の二倍公式で
$2q=(\frac14:-\frac78:1)$である．
Nagell--Lutzの定理により$q$は無限位数である：
有限位数なら$2q$も有理torsionで整数座標を持つはずだが
$\frac14$は整数でない．
$N=(1,3,2,5,7,11)$の六点について
20通りの三点行列式と六点のconic評価行列式が
全て非零であり，72根の共鳴指数も全て非零であることを
\texttt{computations/verify\_resonance\_maximal\_a2.py}
により\emph{有理数の厳密演算}で検算した．
このコードは$A$の表現行列や$d$のrankを
計算するものではなく，Rainsのsheaf存在を
数値実験で置き換えてはいない．
\end{remark}

\begin{proposition}[全次数の最大共鳴]\label{thm:all-degree-maximal}
次数$d=1,\ldots,6$に対し，$p_i=N_iq$（相異なる正整数）なら，全根は非零指数で数値共鳴する．この事実と任意の整数順序を指定したZariski稠密性の完全な証明は主論文\cite{Hearts}にある．
\end{proposition}
\section{最大共鳴での既約根sheaf・zigzag Ext代数・二次障害}
\label{sec:resonance-zigzag}
本節では次数$d=1,\ldots,6$の最大共鳴点について，
「数値的に共鳴するすべての根が既約sheafを与える」
という誤った推論を排除し，実在するsheafのExt代数を
明示的に決定する．
先行研究からの入力は，Rains
\cite[\S4, pp.~33--34]{RainsFiltered}のsimple-root
effectivity判定と
\cite[Definition 9.1, Proposition 7.19, \S11]{RainsBirational}
の既約根sheafに関する結果である．
根の零でないEuler指数の整数表示は
定理\ref{thm:all-degree-maximal}により与えられている．

次数$d=9-r$，$3\le r\le8$とし，
単純根を
\[
\alpha_0=h-e_1-e_2-e_3,\qquad
\alpha_i=e_i-e_{i+1}\quad(1\le i<r)
\]
とする．単純根が実効的で，境界に制限して零の
既約純一次元sheaf $M_i$として実現される
Rains型smooth elliptic heartを固定する．
定理\ref{thm:all-degree-maximal}の
\emph{全根共鳴}の例ではこの仮定が
Rainsのsimple-root存在結果によって保証される．
$\Delta$を単純根のなすDynkin図とする．

\begin{theorem}[実在する既約根はちょうど単純根]
\label{thm:only-simple-irreducible}
$Li^*M=0$を満たす既約純一次元根sheaf $M$が
同じRains heartに存在するとき，
\[
 c_1(M)\in\{\alpha_0,\ldots,\alpha_{r-1}\}
\]
である．逆に各$\alpha_i$は上の$M_i$で実現される．
$q$が非torsionならば，境界非交叉の既約根sheafは
各$\alpha_i$について同型を除き\emph{一つだけ}である．
従って，数値的に共鳴している根の総数と
実際の既約根sheafの同型類数は
\[
\begin{array}{c|rrrrrr}
d&6&5&4&3&2&1\\\hline
\#\text{数値的共鳴根}&8&20&40&72&126&240\\
\#\text{既約根sheaf}&3&4&5&6&7&8
\end{array}
\]
と異なる．
\end{theorem}
\begin{proof}
根系の正根$\gamma$が単純根でないとする．
$\gamma=\sum_{i=0}^{r-1}c_i\alpha_i$，
$c_i\in\mathbb Z_{\ge0}$と書くと
\[
 -2=\gamma^2=\sum_i c_i(\gamma\cdot\alpha_i)
\]
なので，ある$c_i>0$について$\gamma\cdot\alpha_i<0$である．
しかし既存の命題\ref{prop:root-intersection}によれば，
異なる根類の境界非交叉既約根sheafが同時に存在するなら
$\dim\Ext^1(M_\gamma,M_i)=\gamma\cdot\alpha_i\ge0$
でなければならない．従って$\gamma$は不可能．

負根$\gamma$については$-\gamma$が正根であり，
有効な$\alpha_i$の非負整数和なので$-\gamma$は有効である．
もし$M_\gamma$が存在すれば$\gamma$も有効だから，
Rains\cite[Proposition 7.19]{RainsBirational}の
「$D$と$-D$がともに有効なら$D=0$」に反する．
従って単純正根だけが残る．

一意性を証明する．$c_1(M)=c_1(N)=\alpha_i$で
$Li^*M=Li^*N=0$なら，
restriction determinantの式\eqref{eq:resonance}と
$q$の非torsion性から
$\chi(e,M)=\chi(e,N)$である．
従って$[M]=[N]$ in $K_0^{\mathrm{num}}$で，
Rains\cite[Corollary 7.13]{RainsBirational}により
$\chi(M,N)=-\alpha_i^2=2$．
$\theta M\simeq M,\theta N\simeq N$なので，
もし両向きの$\Hom$が零ならSerre双対性で両向きの
$\Ext^2$も零となり，$\chi(M,N)=-\dim\Ext^1(M,N)\le0$
となって矛盾する．従っていずれかの向きに非零射がある．
既約性と純一次元性により，非零射のkernelは零である：
もし非零kernelがあれば同じChern類を持ち，
像は零次元となってtargetの純一次元性に反する．
両者の数値類が同じなので，
Rains\cite[Corollary 7.7]{RainsBirational}から
単射は同型である．
\end{proof}

\begin{theorem}[Dynkin zigzag Ext代数]
\label{thm:zigzag-ext}
$T=\bigoplus_{i=0}^{r-1}M_i$と置く．
$\operatorname{Ext}^{\bullet}_{\mathcal A_{\mathrm R}}(T,T)$は
次の明示的なgraded algebra $Z_\Delta$と同型である．
\begin{enumerate}
\item 次数$0$には各頂点の冪等元$1_i$が一個ずつある．
\item 次数$1$には$\Delta$の各辺$i-j$に対して
両向きの矢$a_{ij}$，$a_{ji}$が各一個ある．
次数$2$には各頂点に一個のloop $z_i$がある．
\item 異なる頂点へ向かう長さ二の道の積は零．
各辺$i-j$の往復積は
$a_{ij}a_{ji}=c_{ij}z_i$，
$a_{ji}a_{ij}=c_{ji}z_j$（$c_{ij},c_{ji}\ne0$）である．
頂点基底と矢をrescaleすれば通常のzigzag関係式にできる．
\item 次数$>2$の積はすべて零．孤立頂点（$A_1$成分）
も次数$2$の$z_i$を持つ．
\end{enumerate}
従って$\#V(\Delta)=r$，$\#E(\Delta)=s$とすると
\[
\dim_k\operatorname{Ext}^{\bullet}(T,T)=2r+2s
=\begin{cases}
8,&d=6,\\
14,18,22,26,30,&d=5,4,3,2,1\text{ の順}.
\end{cases}
\]
これはまず\emph{次数付きYoneda積}の同定である．
dgおよびminimal $A_\infty$模型の形式性については，
以下でLiu--Wangの既知定理を用いて判定する．
\end{theorem}
\begin{proof}
定理\ref{thm:sign}と命題\ref{prop:spherical}によれば
$\Ext^0(M_i,M_i)=\Ext^2(M_i,M_i)=k$，
$\Ext^1(M_i,M_i)=0$．
異なる$i,j$については
命題\ref{prop:root-intersection}より
\[
\Ext^0(M_i,M_j)=\Ext^2(M_i,M_j)=0,\qquad
\dim\Ext^1(M_i,M_j)=\alpha_i\cdot\alpha_j,
\]
これは隣接頂点で$1$，それ以外で$0$である．
高次Extは滑らかなRains heartのhomological dimension $2$
から消える．$\theta\to\mathrm{id}$は全$M_i$上で同型だから，
これらが生成する部分三角圏ではSerre関手が自然に$[2]$となる．
よってSerre traceは
\[
\Ext^1(M_i,M_j)\otimes\Ext^1(M_j,M_i)
\longrightarrow \Ext^2(M_i,M_i)\xrightarrow{\operatorname{tr}} k
\]
を完全な双線型pairingにする．
各辺で往復積が非零となり，
二つの異なる頂点間の$\Ext^2$消滅から
非往復の長さ二の積は零である．
全頂点に次数二の一次元空間が一つずつあるため，
$\Delta$が木（$A_2+A_1$では森）であることを用い
矢とloopの基底をrescaleすれば
各頂点で往復積が同じloopとなる．
以上がzigzag代数の表示である．次数別の次元は
$r,2s,r$なので合計$2r+2s$となる．
\end{proof}

\begin{corollary}[二次Kuranishi障害のpreprojective型表示]
\label{cor:zigzag-quadratic}
$T=\bigoplus_iM_i$の一次変形を
$x=\sum_{i-j}x_{ij}a_{ij}\in\Ext^1(T,T)$
と書く．その二次Kuranishi障害
$m_2(x,x)\in\Ext^2(T,T)$は，適切な非零定数$c_{ij}$を用い
\[
m_2(x,x)=\sum_i
\left(\sum_{j\sim i}c_{ij}x_{ij}x_{ji}\right)z_i
\]
である．$\Delta$が森なので，この二次障害の消滅条件は
\[
\boxed{x_{ij}x_{ji}=0\qquad(\{i,j\}\in E(\Delta))}
\]
と\emph{イデアルとしても}同値である．
従って二次障害方程式で定まるスキーム（gauge商を取る前）は
$2^s$個の座標線型空間の合併で，その各々は次元$s$である．
$d=3$の$E_6$なら$s=5$なので$32$成分となる．
ここまでは一次Yoneda積から得る二次方程式の計算である．
高次障害が追加されないことは次の定理によって従う．
\end{corollary}
\begin{proof}
次数二のYoneda積は前定理の往復積以外は零なので
最初の表示を得る．
任意の葉$i$について次数二の係数方程式は
$c_{ij}x_{ij}x_{ji}=0$ただ一つであり$c_{ij}\ne0$．
従ってその辺積は障害イデアルに属する．
この辺の項を隣の頂点$j$の方程式から取り去って
森の残りのグラフで帰納的に繰り返す．
よって障害イデアルは
$(x_{ij}x_{ji}:\{i,j\}\in E(\Delta))$に等しい．
これは互いに異なる変数対を使う平方自由単項式のイデアルで，
各辺から一方の矢を零にする$2^s$個の座標部分空間を
既約成分とし，各成分では残りの矢が$s$個自由である．
\end{proof}

\begin{theorem}[根配置のdg形式性と形式変形の二次完全性]
\label{thm:zigzag-formal}
基礎体が標数$0$で，$\Delta$が上の
$A_2+A_1,A_4,D_5,E_6,E_7,E_8$のいずれかなら，
$T=\bigoplus_i M_i$の導来自己準同型dg代数は
\[
\boxed{\mathbf R\operatorname{End}_{\mathcal A_{\mathrm R}}(T)
 \ \simeq_{A_\infty}\ Z_\Delta}
\]
として形式的である．従って\emph{固定した}Rainsの
enhanced導来圏の中で$T$の形式的変形を考え，
各直和因子の標識を固定したMaurer--Cartanの
framedパラメータ空間では，
\[
\boxed{
 \widehat{\mathcal O}_{\rm MC}
 \cong
 k[\![x_{ij},x_{ji}:\{i,j\}\in E(\Delta)]\!]/
      (x_{ij}x_{ji}:\{i,j\}\in E(\Delta)).
}
\]
よって高次$A_\infty$障害による追加の切断はなく，
この形式的スキームは$2^{s}$個の
$s$次元formal座標部分空間の合併である
（$s=\#E(\Delta)$）．
特にcubic $E_6$では$32$個の$5$次元formal分枝を持つ．
\end{theorem}
\begin{proof}
Liu--Wang\cite[Theorem~1.1]{LiuWangZigzag}は，
標数$0$における連結ADE Dynkin木の
\emph{path-length graded} zigzag代数が
intrinsically formalであることを示す．
ここで$Z_\Delta$のpath lengthは先のExtの
cohomological degree $0,1,2$に一致している．
Kadeishviliのminimal模型定理により
$\mathbf R\operatorname{End}(T)$の最小$A_\infty$模型は
その$\Ext^\bullet$上に取れるので，
前定理とLiu--Wangを合わせると全ての高次
$m_n$ $(n\ge3)$は$A_\infty$準同型で消せる．

$d=6$の$A_2+A_1$では両連結成分間の
全$\Ext^\bullet$が零なので，minimal模型は
成分ごとに直和へ分離する．
$A_2$にはLiu--Wangを使い，孤立$A_1$の
$\Ext^\bullet=k\oplus kz$（$\deg z=2$，$z^2=0$）
については，strictly unitalな$m_n$ $(n\ge3)$は
正の次数の入力しか持たないと出力次数が
$2n+2-n=n+2>2$になり，次数の理由で零となる．
よってこの場合も形式的である．

Framedな変形のMaurer--Cartan元は
$\Ext^1(T,T)\widehat\otimes\mathfrak m$に属する
$x=\sum x_{ij}a_{ij}$であり，
形式性から方程式は$m_2(x,x)=0$だけで足りる．
前系の葉からの帰納法でその定義イデアルは
$(x_{ij}x_{ji})$そのものとなる．
従って表示した完成環とformal分枝数を得る．
各直和因子の標識を忘れた変形の群oidは
対角の冪等元に対応するgauge群の作用で
さらに商を取るべきであり，ここでは扱わない．
\end{proof}
\begin{remark}[形式性が与えないもの]
この定理はRainsの\emph{固定された}dg圏における
$T$の形式的変形の記述である．
Rains曲面や楕円境界のパラメータを動かす
相対変形，anticanonical heartの全体，
$A=R/(g-1)$の有限次元表現のモジュライ，
および混合根のconnecting map$\delta$の
具体的rankは，この形式性からは決定しない．
\end{remark}

\begin{remark}[導来比較を監査する新しい有限モデル]
特にcubicの最大共鳴で実在する$6$個の単純根sheafの
graded Yoneda algebraは次元$22$の$E_6$ zigzagである．
Rains--anticanonicalの指定した導来切断が完全忠実かを調べるには，
このgraded algebraの積だけでなく，
$\mathbf R\operatorname{Hom}(e,T)$の$A_\infty$-module構造，
とりわけ正負隣接根間の
$\operatorname{Ext}^1$から誘導される
$V_\alpha\to V_\beta$のconnecting morphismを決定する必要がある．
ここでの形式性は上の定理で確立したが，
このことだけでは$\delta$のrankは決まらない．
具体的な$A$加群の作用行列とnatural functorの
$A_\infty$-module構造の計算が独立に必要である．
\end{remark}

\section{非torsion境界外圏と混合\texorpdfstring{$A_2$}{A2}の完全判定}
\label{sec:mixed-a2-category}
本節は自然な導来切断の完全忠実性を仮定せず，混合根のconnecting mapを決定する．
Rains\cite[Corollaries 7.7, 7.13, 7.15, Propositions 7.6, 7.32,
Corollary 10.12, Proposition 10.13]{RainsBirational}と既存の符号別Extを入力とする．
$H=-K_X$はnef，$H^2=d>0$，$q$は非torsion，構造line bundle $e$はexceptionalとする．
$\Ext_X$は常にambientのRains heartで計算する．

\subsection{境界外圏の有限長性}
\begin{lemma}[零次元の境界外商の排除]
\label{lem:offboundary-no-points}
$Li^*M=0$なら，$M$の零次元商は零である．従って境界非交叉の
既約純一次元sheafはRains heartの単純対象である．
\end{lemma}
\begin{proof}
零次元商$M\twoheadrightarrow Z$には右完全性から$i^*Z=0$が従い，
$\theta Z\to Z$は全射となる．両者の数値類は等しいので
Rains Corollary 7.7により同型，従って$Li^*Z=0$である．
一方$[Z]=\ell[\mathrm{pt}]$，$\ell=\chi(e,Z)\ge0$で，
境界restrictionのdeterminantは$q^\ell$となる
\cite[Proposition 7.6, Proposition 7.24直前]{RainsBirational}．
非torsion性から$\ell=0$，ゆえに$Z=0$．既約純一次元sheafの
非零部分対象の商は零次元なので単純性が従う．
これは非torsion性と境界非交叉を使う証明であり，既約性だけから単純性を推論してはいない．
\end{proof}

\begin{proposition}[境界外Serre部分圏]
\label{prop:offboundary-serre-length}
$\mathcal Z=\{M:\dim M\le1,\ Li^*M=0\}$は有限長Serre部分圏であり，
その単純対象は2-sphericalな境界非交叉根sheafである．
最大共鳴で全単純根sheafが実在する場合，単純対象は
定理\ref{thm:only-simple-irreducible}の$r=9-d$個に尽きる．
\end{proposition}
\begin{proof}
$0\to N\to M\to Q\to0$，$M\in\mathcal Z$とする．
nef性と加法性から$0=Hc_1(M)=Hc_1(N)+Hc_1(Q)$の両項は零．
$i^*Q=0$なので$\theta Q\to Q$は全射であり，rank零と$K_Xc_1(Q)=0$から
両者の数値類は同じ\cite[Corollary 7.15]{RainsBirational}．
Corollary 7.7から同型となり，$Li^*Q=0$．自然性から$\theta N\to N$も同型．
拡大閉性は境界triangleによる．
非零零次元対象がないため，真の降鎖の逐次商はすべて一次元である．
Rains Proposition 7.32のfinitely partitive性とnoetherian性から有限長となる．
単純対象$S$の$D=c_1(S)\ne0$は$HD=0$を満たす．
$H^\perp$の負定値性，Riemann--Rochの偶奇，$\theta S=S$とSerre双対性より
\[
0<-D^2=\chi(S,S)=2-\dim\Ext_X^1(S,S)\le2
\]
であり，$D^2$は偶数だから$D^2=-2$，$\Ext_X^1(S,S)=0$．
最後は既存の単純根の排他性と一意性を使う．
\end{proof}

\subsection{任意の正負指数に対する最大rank}
隣接する境界非交叉既約根sheafを$S_+,S_-$とし，
\[
\alpha=c_1(S_+),\quad\beta=c_1(S_-),\quad\alpha\beta=1,
\qquad\chi(e,S_+)=p>0,\quad\chi(e,S_-)=-n<0
\]
とする．$\gamma=\alpha+\beta$は根である．非分裂拡大を
\begin{equation}\label{eq:a2-two-extensions}
0\to S_-\to P\to S_+\to0,\qquad
0\to S_+\to Q\to S_-\to0
\end{equation}
と書く．両向きのExt$^1$は一次元なので各中間項は同型を除き一意である．

\begin{theorem}[任意の正負指数のconnecting rank]
\label{thm:a2-delta-rank}
$P$のconnecting map$\delta:k^p\to k^n$は
\[
\boxed{\operatorname{rank}\delta=\min(p,n),\qquad
 (\dim V^0(P),\dim V^1(P))=((p-n)_+,(n-p)_+)}
\]
を満たす．ここで$x_+=\max(x,0)$．
$p>n$なら非零$e\to P$は全射で，核$L$はrank-one torsion-free，
\[
c_1(L)=-\gamma,\quad\chi(e,L)=1-p+n,\quad
\chi(L,L)=3-2(p-n).
\]
$L$がline bundleとなることと$p-n=1$は同値である．
特に$p=2,n=1$ではrankは1，$(h^0,h^1)=(1,0)$である．
いずれにも導来完全忠実性を仮定しない．
\end{theorem}
\begin{proof}
符号別Extから$\Hom(e,P)=\ker\delta$，$\Ext^1(e,P)=\operatorname{coker}\delta$．
非零$f:e\to P$があるとする．$\Hom(e,S_-)=0$なので$P\to S_+$との合成は非零．
$S_+$の単純性から全射である．像$I$と$S_-$の交わりは$0$か$S_-$であり，
前者なら拡大を分裂させてしまう．従って$I=P$，すなわち$f$は全射．
その核$L$はrank-one torsion-freeで$c_1(L)=-\gamma$である．
Serre双対性とRains Corollary 10.12より
\[
\Ext^2(e,L)=\Hom(L,\theta e)^*=0:
\quad H(c_1(\theta e)-c_1(L))=H(K_X+\gamma)=-d<0.
\]
ここでCorollary 10.12はrank-one torsion-free同士に適用でき，
$L$をline bundleと仮定する必要はない．
$e$のexceptionalityと$0\to L\to e\to P\to0$から$\Ext^1(e,P)=0$となる．

$p>n$なら$\ker\delta\ne0$なので上の議論が適用でき，$\delta$は全射．
$p\le n$なら$\Hom(e,P)\ne0$とすると$h^1(P)=0$になり，
$0\ge p-n=\chi(e,P)=h^0(P)>0$となって矛盾する．
従ってこの場合は$\delta$が単射．これで最大rankを得る．
$p>n$の場合，Euler加法性とRiemann--Rochから表示した$L$の値を得る．
Rains Proposition 10.13の等号条件$\chi(L,L)=1$により，
line bundle条件は$p-n=1$と同値である．
\end{proof}
\begin{remark}[指数の釣合いと消失]
$p=n$で上の二sheafが存在すれば，$P\ne0$なのに$\Phi(P)=0$となる．
しかしgenuine smooth classical avatarでは$\gamma$が零指数の共鳴根となるため，
この釣合いの場合自体が起こらない．一般の数値条件からこのsheafの実在を主張しない．
\end{remark}

\begin{corollary}[二段直和拡大]
\label{cor:a2-rank-matrix-exact}
$0\to S_-^{\oplus b}\to M\to S_+^{\oplus a}\to0$の拡大類を
$T\in\operatorname{Mat}_{b\times a}(k)$，$t=\operatorname{rank}T$，$r=\min(p,n)$とすると
\[
\operatorname{rank}\delta_M=rt,\qquad
\dim V^0(M)=pa-rt,\qquad\dim V^1(M)=nb-rt.
\]
$p=2,n=1$では$(2a-t,b-t)$となり，$a\ge1$ならEuler標数が零でも消失しない．
\end{corollary}
\begin{proof}
Yoneda積の自然性から$\delta_M=\delta\otimes T$であり，rankは積となる．
六項完全列を適用する．
\end{proof}

\subsection{有限長部分圏の全対象と局所torsion pair}
\begin{theorem}[四つの直既約対象]
\label{thm:a2-serre-category}
$\mathcal C=\langle S_+,S_-\rangle_{\rm Serre}$は
$\operatorname{mod}_{\rm fd}\Pi(A_2)$と同値であり，直既約は$S_+,S_-,P,Q$だけである．
$P\oplus Q$はprojective generatorで，$\Pi(A_2)$は二頂点・両向き二矢と
両方の長さ二の道を零とする関係を持つ次元4の代数である．
\end{theorem}
\begin{proof}
非分裂性と単純性から，順序$(S_+,S_-,P,Q)$で
\[
(\dim\Hom(U,V))_{U,V}=
\begin{pmatrix}1&0&0&1\\0&1&1&0\\1&0&1&1\\0&1&1&1\end{pmatrix}.
\]
例えば$\Hom(S_+,P)=\Hom(P,S_-)=0$は，connecting map
$k\to\Ext_X^1(S_+,S_-)=k$が非零拡大類を送ることから従う．
$\Hom(P,Q)=k$は$P\twoheadrightarrow S_+\hookrightarrow Q$で生成され，
逆方向も同様で，両合成は零である．
$\theta$が全対象で自然に恒等なので
\[
\dim\Ext_X^1(U,V)=\dim\Hom(U,V)+\dim\Hom(V,U)+c_1(U)c_1(V).
\]
上の表を代入すると，$S_+,S_-$間の両向きの一次元Ext以外は全て零となる．
とくに$P,Q$から両単純へのExt$^1$は零．長さによる帰納法から
$P,Q$は$\mathcal C$でprojectiveで，単純への全射を持つのでgeneratorとなる．
そのEnd代数のHomと合成が上記の表示を与える．
表現を$x:U_+\to U_-$，$y:U_-\to U_+$，$xy=yx=0$で書くと，
$\operatorname{im}y\subset\ker x$，$\operatorname{im}x\subset\ker y$に適合する
基底により，一方の矢のみ非零の一次元対と孤立一次元空間へ分解される．
これが四つの直既約である．
\end{proof}

\begin{corollary}[全cohomologyと指数に応じた局所torsion pair]
\label{cor:a2-torsion-pair}
$\mathcal C$の直既約の$(h^0,h^1)$は
\[
S_+:(p,0),\quad S_-:(0,n),\quad P:((p-n)_+,(n-p)_+),\quad Q:(p,n).
\]
一般のquiver表現で$a=\dim U_+$，$b=\dim U_-$，$t=\operatorname{rank}x$なら
\[
h^0=pa-\min(p,n)t,\qquad h^1=nb-\min(p,n)t.
\]
$p>n$のとき$(\ker h^1,\ker h^0)=(\operatorname{add}(S_+,P),\operatorname{add}(S_-))$，
$p<n$のとき$(\operatorname{add}(S_+),\operatorname{add}(S_-,P))$はtorsion pairである．
$p\ne n$なら$\Phi$は非零対象を消さない．$p=n$の場合の核は$\operatorname{add}(P)$で，
二つのkernel部分圏はtorsion pairにならない．
\end{corollary}
\begin{proof}
$P$は最大rank定理，$Q$は逆向きの六項完全列から求まる．
四つの直既約の重複度は順に
$a-\operatorname{rank}x-\operatorname{rank}y$，
$b-\operatorname{rank}x-\operatorname{rank}y$，
$\operatorname{rank}x$，$\operatorname{rank}y$であり，直和して公式を得る．
Hom表がtorsion pairの直交性を与え，$Q$の短完全列がtorsion分解を与える．
他の直既約は片方の部分圏に属する．非消失および核の記述もこの表から従う．
\end{proof}
\begin{remark}
$\chi(e,P)=\chi(e,Q)=p-n$だが，$Q$には二つのcohomological degreeが残る．
非既約対象ではEuler指数の符号だけからheart内の位置を決定できない．
また$\Ext_{\mathcal C}^{2j}(S_+,S_+)=k$ $(j\ge0)$とambientの
$\Ext_X^{>2}=0$は異なる．従って$D^b(\mathcal C)\to D^b(\mathcal A_q)$は
完全忠実でなく，このabel分類をambient三角圏の同定に読み替えない．
\end{remark}

\subsection{二次元表現の検査条件と導来比較への接続}
$p=2,n=1$に戻り，$V_+=V^0(S_+)$，$V_-=V^1(S_-)$，$W=V^0(P)$とする．
\begin{proposition}[単一Homテストの正確な範囲]
\label{prop:a2-scalar-test}
$0\to W\to V_+\xrightarrow{\delta}V_-\to0$は$A$加群の短完全列で，両端は一次元．
$\dim\Hom_A(V_+,V_-)$は$V_+\simeq V_-^{\oplus2}$の場合のみ2，それ以外は1である．
自然写像$\Ext_X^1(S_+,S_-)\to\Hom_A(V_+,V_-)$は常に単射である．
\end{proposition}
\begin{proof}
六項完全列と最大rankから短完全列を得る．準同型空間は
二次元の$\Hom_k(V_+,V_-)$の非零部分空間．独立な二準同型があれば
それらを並べて$V_+\simeq V_-^2$を得る．拡大類の生成元が非零$\delta$に写るので単射性も従う．
\end{proof}

右加群を行ベクトルで表し，$W$を第一座標に取ると
\[
\rho(a)=\begin{pmatrix}\chi_W(a)&0\\\eta(a)&\chi_-(a)\end{pmatrix},\qquad
\eta(ab)=\eta(a)\chi_W(b)+\chi_-(a)\eta(b).
\]
\begin{proposition}[導来完全忠実性の必要条件]
\label{prop:a2-character-obstructions}
$\Phi$が二根の生成する三角圏上で完全忠実なら，$W\not\simeq V_-$で，
上の加群拡大は非分裂である．従って作用像は全三角行列代数である．
\end{proposition}
\begin{proof}
$\Ext_X^1(P,S_-)=0$から完全忠実性の下で$\Hom_A(W,V_-)=0$．
また非零包含$S_-\hookrightarrow P$は加群拡大の非零類へ移る．
異なるcharacterと非分裂性により二つの対角冪等元と非零非対角成分を得る．
\end{proof}
\begin{remark}
同じcharacterのJordan型は単一Hom次元1のテストには合格しても上の必要条件に反する．
異なるcharacterでも分裂型なら非零sheaf射を消す．対象の非消失と射の忠実性は別である．
ただし本稿のgenuine楕円gluingでは，次節の比較定理と
定理\ref{thm:resonance-affine-triangular}により両方の必要条件が実際に成立する．
残るのは元の$A$の生成元に対するcharacterと非対角係数の具体的な式であって，
この範囲での分離・非分裂性の真偽ではない．
\end{remark}



\section{零指数とgenuine性の境界}\label{sec:resonance-zero-summary}
\begin{proposition}[genuine avatarにおける零指数の排除]\label{prop:zero-index}
滑らかなgenuine del Pezzo $S$と滑らかな反標準曲線$E$について，任意の根$\alpha$は$\mathcal O_E(\alpha)\ne\mathcal O_E$を満たす．従って非torsionな$q$に関して零指数根共鳴はない．
\end{proposition}
\begin{proof}
$-K_S$のamplenessと根のRiemann--Rochから$H^0(S,\pm\alpha)=0$であり，制限完全列から$H^0(E,\mathcal O_E(\alpha))=0$を得る．詳細は主論文\cite{Hearts}を参照する．
\end{proof}
零指数壁の古典的収縮・非可換弱gluingのorbit核とdg商は主論文で扱う．

\section{再現と今後の重点}
符号別Euler計算と根格子は\texttt{computations/verify\_resonance\_hearts.py}，
拡大の相殺と根の交叉は\texttt{computations/verify\_resonance\_extensions.py}，
全根共鳴の有理楕円点・全次数の根数・整数指数は
\texttt{computations/verify\_resonance\_maximal\_a2.py}，
既約根の排除条件とzigzag分枝数は\texttt{computations/verify\_resonance\_zigzag.py}で検算する．
混合$A_2$のHom/Euler/Ext，任意正負指数の両chamber，四次元商の行列模型，
および局所Weyl反射は\texttt{computations/verify\_resonance\_mixed\_a2.py}で検算する．
これらはsheaf存在や導来比較を計算で証明するものではない．

現在の具体的な未解決問題は，元の楕円データから$A=R/(g-1)$の生成元を固定し，
商$A\twoheadrightarrow\Pi(A_2)$の各係数を明示することである．
characterの分離，二次元拡大の非分裂性，および指定gluingでの全heartのtorsion pairは
もはや未決定として扱わない．一方，全$E_6$普遍有限商での楕円生成元の具体的作用行列，独立のavatar間での指定同定，相対family比較と退化境界は引き続き課題である．
実行済み検算・組版・文献照合の正確な範囲は
\texttt{research/resonance\_hearts\_validation.json}に記す．



\part{三次曲面の一頂点cornerと六点配置}
\section{対象・grading・既知入力}
基礎体は標数0の代数閉体とする．頂点を
$A_a,B_b,C_c$ ($0\le a,b,c\le2$)，矢を
\[
 x_{ab}:A_a\to B_b,\qquad y_{bc}:B_b\to C_c,\qquad z_{ca}:C_c\to A_a
\]
とする．道の積は右から左への合成で，potentialとJacobi環は
\[
 W(w)=\sum_{a,b,c}w_{abc}[z_{ca}y_{bc}x_{ab}],\qquad
 B(w)=\Jac(Q,W(w)),\qquad e=e_{A_0}
\]
である．このquiverとpotentialはAOU\cite[\S4]{AOU}のcubic型である．
$w$を動かす係数空間は$\mathbb A^{27}$．全ての矢の次数を1とする
\emph{道長grading}を$\ell$と書く．全関係は道長2である．
一方，\emph{cut grading}は
\[
 \deg x=\deg y=0,\qquad\deg z=1
\]
である．反標準環と$qgr$の比較には後者を用い，$B^{\mathrm{cut}}$と書く．
両gradingで環の非次数付き構造は同じである．各cut次数の道長には有限の上界があるので，
$B^{\mathrm{cut}}$も局所有限で，次数0部分は有限次元である．

AOU\cite[(3.19), (4.9), (4.13)]{AOU}の標準的なhelix同定を使うAS領域では，
\[
 \cA_{\mathrm{AS}}\simeq\qgr B^{\mathrm{cut}},\qquad
 O=\pi(eB),\qquad \tau=\mathbb S^{-1}[2]=(1)_{\mathrm{cut}}.
\]
ここではfoldingがこの同定を与える標準的なmarked AS対象を扱う．
AS次元3のprojectiveの飽和から，cut次数付き同型
\begin{equation}\label{eq:orbit}
 eB^{\mathrm{cut}}e\simeq R(O,\tau)
 =\bigoplus_{n\ge0}\Hom_{\cA_{\mathrm{AS}}}(O,\tau^nO)
\end{equation}
を得る．飽和の理由は，AS条件により単純対象からprojectiveへの
$\Ext^0,\Ext^1$が消え，有限長加群へ組成列で拡張できることにある．
$qgr$のHomを定める有限長の切断を取ってもHomは変わらない．
右加群のorbit積は$ab=\tau^{\deg b}(a)\circ b$とする．

\begin{remark}
道長gradedの$\qgr B$を，そのまま一つのAS heartと同一視しない．
blockと道長の合同類による分解があるためである．
cornerの道長は3の倍数であり，反標準次数はその3分の1である．
cut gradingでfoldingすればこの曖昧さはない．
一般の係数でAS正則性やSerre同定を新たに証明したのではなく，
上記の標準AS領域に限定して適用する．
\end{remark}

\section{有限次元corner商と開性}
\begin{lemma}\label{lem:finite}
$B$は有限quiverの道長斉次関係による商で，$e$は頂点の和とする．
$C=B/BeB$に対して$C_N=0$なら，$C_n=0$ for $n\ge N$である．
このとき$Be$は右$eBe$加群，$eB$は左$eBe$加群として有限生成である．
\end{lemma}
\begin{proof}
$C$は矢で生成されるので，$C_N=0$の後は乗法によって全成分が零になる．
従って道長$N$の道は，$e$を通る同じ長さの道の線形結合に書ける．
$Be$の長い道の終端側$N$個をこの形で置き換えると，
右側に$eBe$の元を取り出し，左側を長さ$\le N$の$Be$の道にできる．
$e$で切る位置が端にある場合も，短い側の長さを$N$まで許せばよい．
残る場合は長さの帰納法を使う．よって長さ$\le N$の有限個の道が右加群を生成する．
逆向きの道について同じ議論を行えば$eB$の左有限生成性を得る．
この証明はcut gradingにも斉次的である．
\end{proof}

\begin{proposition}\label{prop:corner}
$B^{\mathrm{cut}}$が右noetherianであるとする．$C=B/BeB$が有限次元なら，
自然な関手
\begin{equation}\label{eq:equiv}
 \qgr B^{\mathrm{cut}}\longrightarrow\qgr(eB^{\mathrm{cut}}e),\qquad M\longmapsto Me
\end{equation}
はabel同値である．逆に，このrestrictionが$qgr$で同値なら$C$は有限次元である．
\end{proposition}
\begin{proof}
$R=eBe$は右noetherianである．実際，$R$の右イデアル$I$に対し
$(IB)e=I$なので，右イデアルの増大列は$B$のnoetherian性で止まる．
補題\ref{lem:finite}より$Me$は有限生成$R$加群となる．
restrictionは完全で，有限次元加群を保つ．
左随伴$L(N)=N\otimes_R eB$は有限生成加群を有限生成加群に送り，
$eB$の左有限生成性から有限次元加群も保つ．
unitは$(N\otimes_R eB)e\simeq N$であり，counitの核と余核は$e$で消える．
核は$B$の右noetherian性で有限生成である．
両者は有限次元環$C$の有限生成加群なので有限次元である．
$L$が短完全列の左端に作る核も同じ理由で有限次元となるから，
$qgr$上で$L$は完全となり，unit/counitが同値を与える．
逆に$\pi C$はrestrictionで零に送られる．同値なら$\pi C=0$であり，
$C$は有限生成$B$加群なので有限次元である．
\end{proof}

$V_N$を$A_0$を除くquiverの道長$N$の道を基底とするベクトル空間とする．
全ての$u(\partial_aW)v$，$\ell(u)+\ell(v)=N-2$を列挙した係数行列を$M_N(w)$とする．
その成分は$w$の線形式であり，$C_N$はその行空間による$V_N$の商である．
\begin{theorem}\label{thm:open}
\[
 U_N=\{w:\operatorname{rank}M_N(w)=\dim V_N\},\qquad
 U=\bigcup_{N\ge2}U_N
\]
は係数空間の開集合であり，$U$は$\dim C(w)<\infty$と同値である．
さらに，ある一様な整数$N_0$が存在し$U=U_{N_0}$となる．
$U$上のnoetherian AS対象では
\[
 \cA_{\mathrm{AS}}\simeq\qgr R(O,\tau)
\]
であり，この同値は構造対象と反標準shiftを保つ．
従ってこの範囲では導来同値だけでなくheart自体が一致する．
\end{theorem}
\begin{proof}
最大minorの非消滅が$U_N$を定め，補題\ref{lem:finite}から
$U_N\subset U_{N+1}$となる．有限次元性はある$N$での消滅と同値．
係数空間$\mathbb A^{27}$はnoetherianなので開集合$U$は準コンパクトであり，
この昇鎖はある$N_0$で$U$を覆う．$N_0=7$とは主張しない．
命題\ref{prop:corner}と式\eqref{eq:orbit}を適用する．
$\pi(eB)$の像は$\pi R$であり，restrictionはcut shiftと可換する．
\end{proof}

\begin{remark}
これは\emph{指定したcorner restriction}の一致開集合である．
$U$の外で別のabel同値や導来同値が存在しないとは結論しない．
全potentialがAS正則，noetherianであるとも，AS locus全体が有限型代数stackであるとも
本定理では仮定・主張しない．射の定義域を同じ有限係数chartで固定した場合の結論である．
\end{remark}


\section{係数トーラスの四変数化と無限次元cornerの障害}
\label{sec:torus}
$C(w)=B(w)/B(w)e_{A_0}B(w)$を考える．$A_0$に接する矢は$C(w)$で消えるので，
この商は$w_{abc}$ ($a=1,2$)の18係数だけに依存する．
それらが全て非零の範囲で
\[
q_{bc}=\frac{w_{2bc}}{w_{1bc}},\qquad
\lambda_{ij}=\frac{q_{ij}q_{00}}{q_{i0}q_{0j}}
\quad(1\le i,j\le2)
\]
と定める．
\begin{theorem}[四交比と九つの必要条件]\label{thm:torus}
\begin{enumerate}
\item 個々の矢の非零rescalingを除いた18係数の軌道は，
四つの$\lambda_{ij}\in k^\times$で完全に分類される．
\item $C(w)$が有限次元ならば18係数は全て非零であり，さらに
\begin{equation}\label{eq:minor}
q_{bc}q_{b'c'}-q_{bc'}q_{b'c}\ne0
\quad(b<b',\ c<c')
\end{equation}
の九条件が成立する．
\item $C(w)$が有限次元になるlocusは，
四次元係数トーラス上の非空の真の開集合である．
以前の可換六点witnessでは
\[
(\lambda_{ij})=
\begin{pmatrix}-1/4&1/16\\1/6&-1/4\end{pmatrix},\qquad C_7=0.
\]
\end{enumerate}
\end{theorem}
\begin{proof}
まず$y_{bc}$をrescaleして全$w_{1bc}$を1にできる．
残りの$x_{2b}/x_{1b}$および$z_{c2}/z_{c1}$のrescalingにより，
$q_{bc}$は任意の行因子・列因子を掛けられる．
$q_{00}=q_{b0}=q_{0c}=1$と正規化すれば
\[
Q=\begin{pmatrix}1&1&1\\1&\lambda_{11}&\lambda_{12}\\
1&\lambda_{21}&\lambda_{22}\end{pmatrix}
\]
だけが残る．逆に四交比はrescalingで不変なので，軌道を完全に記述する．

もし$w_{abc}=0$ ($a=1,2$)ならば，
$x_{ab},y_{bc},z_{ca}$だけを残し他の矢を零にすると
全てのJacobi関係が消える．各矢を$t$倍の行列単位に送り，
残る長さ3の閉路は非零な$t^3$倍の対角行列単位になる．
その全ての冪が非零なので$C(w)$は無限次元である．

全18係数が非零として式\eqref{eq:minor}の一つが零とする．
選んだ二行の添字集合$I$と二列の添字集合$J$で
$q_{bc}=r_bs_c$ ($b\in I,c\in J$)と分解できる．
行和・列和が零で全成分非零の
\[
T=\begin{pmatrix}1&-1\\-1&1\end{pmatrix}
\]
を$I\times J$に置く．$b\in I,c\in J$上で
\[
x_{1b}=1,\quad x_{2b}=-r_b^{-1},\quad
z_{c1}=1,\quad z_{c2}=s_c^{-1},\quad
y_{bc}=\frac{T_{bc}}{w_{1bc}}
\]
とし，残りの矢は零とする．Jacobi関係の左辺は
$T$の行和・列和，または
$w_{1bc}-w_{2bc}/(r_bs_c)$の倍数だからすべて零である．
各非零矢を$t$倍の行列単位に置き換えると，
支持に含まれる閉路が非零な$t^3$倍の対角行列単位となり，
$C(w)$は無限次元である．

有限次元性のlocusは前節の定理により開である．
四交比による正規化はこの条件を四次元トーラス上へ降ろす．
六点witnessの係数を代入すると表示した交比を得て，
既知の$C_7=0$より非空．
式\eqref{eq:minor}の零集合を含まないので真の開集合である．
\end{proof}

\begin{remark}[九つの明示的な禁止divisor]
四交比で式\eqref{eq:minor}は
\[
\begin{gathered}
\lambda_{11},\lambda_{12},\lambda_{21},\lambda_{22}\ne1,\\
\lambda_{11}\ne\lambda_{12},\ 
\lambda_{21}\ne\lambda_{22},\
\lambda_{11}\ne\lambda_{21},\
\lambda_{12}\ne\lambda_{22},\\
\lambda_{11}\lambda_{22}\ne\lambda_{12}\lambda_{21}
\end{gathered}
\]
に等しい．これら九つだけでは不十分であり，二つの行列式障害を加えた十一条件の必要十分性を定理\ref{thm:general-position-iff}で示す．AS locusとの交叉を調べることなしに
heartの非一致や導来非同値を結論しない．
\end{remark}
\section{行列式で検出する追加の無限次元障害}
\label{sec:ranktwo-obstructions}
前節の九つの二次小行列式条件だけでは，corner の有限次元性に
十分ではない．$w_{1bc}\ne0$ を仮定し，矢 $y_{bc}$ の rescaling により
$w_{1bc}=1$ として，$w_{2bc}=q_{bc}$ を記す．
\emph{成分ごとの逆数}からなる行列を
\[
 Q=(q_{bc})_{0\le b,c\le2},\qquad
 Q^{\langle-1\rangle}=(q_{bc}^{-1})_{0\le b,c\le2}
\]
と書く．後者は通常の逆行列 $Q^{-1}$ ではない．

\begin{proposition}[二つの行列式障害]\label{prop:two-determinants}
全ての $q_{bc}\in k^\times$ とする．もし
\[
 \det Q=0\quad\text{または}\quad
 \det Q^{\langle-1\rangle}=0,
\]
ならば $C(w)=B(w)/B(w)e_{A_0}B(w)$ は無限次元である．
従って $C(w)$ の有限次元性には，前節の九つの小行列式の
非零性に加えて，上の二つの行列式の非零性が必要である．
これら二つは九条件からは従わない．
\end{proposition}
\begin{proof}
まず $\det Q=0$ を仮定する．$Q$ を
\[
 Q=UV,\quad U\in M_{3\times2}(k),\quad V\in M_{2\times3}(k)
\]
と分解できる．$0\ne h\in\ker U^{\mathsf t}$ および
$0\ne k\in\ker V$ を選ぶ．それぞれが非零なので，
$h_bk_c\ne0$ となる $b,c$ が存在する．
$u_b$ を $U$ の第 $b$ 行，$v_c$ を $V$ の第 $c$ 列とする．

$k[t]$ 上の quiver 表現を
\[
 A_1=k[t]^2,\quad A_2=k[t],\quad B_b=C_c=k[t]
\]
とし，矢の写像を
\begin{equation}\label{eq:ranktwo-representation}
 x_{1b}=t u_b,\quad z_{c1}=-t v_c,\qquad
 x_{2b}=z_{c2}=t,\quad y_{bc}=t h_bk_c
\end{equation}
で定める．$x_{1b}$ は行ベクトル，$z_{c1}$ は列ベクトルである．
$y_{bc}$ に関する Jacobi 関係は
\[
 x_{1b}z_{c1}+q_{bc}x_{2b}z_{c2}
 =t^2(-u_bv_c+q_{bc})=0.
\]
$x_{1b},z_{c1}$ に関する関係はそれぞれ
$Vk=0$ と $h^{\mathsf t}U=0$ に帰着する．
$x_{2b},z_{c2}$ に関する関係は $Qk=0$ と
$h^{\mathsf t}Q=0$ に帰着する．
従って \eqref{eq:ranktwo-representation} は全ての Jacobi 関係を満たす．
しかも $A_2$ に基点を持つ三角閉路
\[
 z_{c2}y_{bc}x_{2b}:A_2\longrightarrow A_2
\]
の像は $t^3h_bk_c\ne0$ であり，その任意の冪の像も非零である．
よって $C(w)$ に非零元が無限に存在し，無限次元となる．
ここで $k[t]$ は表示の標的であり，有限次元の quiver 表現を
仮定していない．

次に $\det Q^{\langle-1\rangle}=0$ とする．
頂点 $A_1,A_2$ および対応する $x,z$ を交換すると，
残る二系列の potential 係数は $(q_{bc},1)$ となる．
各 $y_{bc}$ を非零スカラーで rescale して第一系列を1に戻せば，
新しい比行列は $Q^{\langle-1\rangle}$ である．
この操作は頂点を入れ替えた graded quiver の同型なので，
前半を適用できる．
\end{proof}

\begin{corollary}[四交比座標での十一の必要条件]
\label{cor:eleven}
\[
 Q\sim
 \begin{pmatrix}1&1&1\\1&a&b\\1&c&d\end{pmatrix}
 \qquad(a,b,c,d\in k^\times)
\]
と規格化する．$C(w)$ が有限次元ならば，前節の九つの二次小行列式が
全て非零であり，さらに
\begin{equation}\label{eq:two-extra-polynomials}
 (a-1)(d-1)-(b-1)(c-1)\ne0,\qquad
 bc(1-a)(1-d)-ad(1-b)(1-c)\ne0
\end{equation}
が必要である．
\end{corollary}
\begin{proof}
最初の式は $\det Q$ である．第二の式は
$abcd\det Q^{\langle-1\rangle}$ に等しい．
前節の定理と命題\ref{prop:two-determinants}を適用する．
\end{proof}

\begin{remark}[九つの小行列式だけでは不十分]
\label{ex:ranktwo}
\[
 Q_0=\begin{pmatrix}3&4&5\\5&6&7\\7&8&9\end{pmatrix},\quad
 Q_1=2520Q_0^{\langle-1\rangle}
 =\begin{pmatrix}840&630&504\\504&420&360\\360&315&280\end{pmatrix}.
\]
両方とも9つの $2\times2$ 小行列式は全て非零である．
一方，
\[
 \det Q_0=0,\quad
 \det Q_0^{\langle-1\rangle}=\frac1{198450},\qquad
 \det Q_1=80640,\quad
 \det Q_1^{\langle-1\rangle}=0.
\]
従って，両方とも $C$ は無限次元である．これらの厳密な有理数演算と
式\eqref{eq:ranktwo-representation}の全 Jacobi 関係の検査は
\texttt{computations/verify\_cubic\_corner\_rank2.py}で再現できる．
\end{remark}

\begin{remark}[有限体での探索と限界]
十一の非零条件の十分性は，次節の定理\ref{thm:general-position-iff}で六点配置との同定によって証明する．
規格化した四交比トーラスを $\mathbb F_7$ 上で全探索すると，
十一の条件を満たす80点すべてで $C_7=0$ になる．
これは有限体の有限個の点に関する独立検算である．標数0の十分性は次節の幾何学的証明に依存し，この有限体計算に依存しない．
同じスクリプトのオプション \texttt{--scan-f7} で再現できる．
\end{remark}

\section{四交比 corner の完全判定と六点配置による十分性}
\label{sec:six-point-classification}
本節では，前節までに証明した十一の必要条件が十分でもあることを示す．
その鍵は，四交比トーラスを\emph{標識付き可換三次曲面の六点配置空間}と同定することである．
以下，基礎体は標数0の代数閉体とし，十八の係数$w_{1bc},w_{2bc}$はすべて非零とする．

\begin{theorem}[corner有限性の必要十分条件]
\label{thm:general-position-iff}
$B(w)=\operatorname{Jac}(\widetilde Q,W(w))$，$e=e_{A_0}$，
$C(w)=B(w)/B(w)eB(w)$と置く．
$Q=(q_{bc})$，$q_{bc}=w_{2bc}/w_{1bc}$とすると，次は同値である．
\begin{enumerate}
\item $\dim_kC(w)<\infty$．
\item $Q$の九つの$2\times2$小行列式がすべて非零であり，
$\det Q\ne0$かつ$\det Q^{\langle-1\rangle}\ne0$である．
\item $P=(Q^{\langle-1\rangle})^{\mathsf t}$の三つの列$p_1,p_2,p_3$と，
座標点$p_4=(1:0:0),p_5=(0:1:0),p_6=(0:0:1)$が，
$\mathbb P^2$の一般位置の六点をなす．
\end{enumerate}
従って，四交比トーラスにおける有限corner locusは，
\emph{正確に}前節の十一のdivisorの補集合である．
この開集合は，六点の順序を保つ標識付き可換三次曲面のモジュライ
$M_{\mathrm{cubic}}$と自然に同型になる．
\end{theorem}
\begin{proof}
(i)$\Rightarrow$(ii)は定理\ref{thm:torus}と
命題\ref{prop:two-determinants}による．
(ii)$\Leftrightarrow$(iii)を示す．$P$の各成分は$q_{bc}^{-1}\ne0$であり，
その$2\times2$小行列式は，添字と非零係数を除けば
$Q$の対応する小行列式と等しい．さらに
$\det P=\det Q^{\langle-1\rangle}$である．従って(ii)のうち
九小行列式と後者の行列式は，順に「可動点二つと座標点一つが非共線」
「可動点三つが非共線」を表す．非零成分は「可動点一つと座標点二つが非共線」を表す．
座標点三つはもとより非共線である．

残る条件は六点が一つの二次曲線上にないことである．
座標三点を通るconicの方程式は
$\alpha X_1X_2+\beta X_2X_3+\gamma X_3X_1=0$
であり，可動点$p_i=(P_{0i}:P_{1i}:P_{2i})$への評価行列は
\[
 M_{i\bullet}=(P_{0i}P_{1i},P_{1i}P_{2i},P_{2i}P_{0i})
 =(P_{0i}P_{1i}P_{2i})
 (P_{2i}^{-1},P_{0i}^{-1},P_{1i}^{-1}).
\]
すべての因子が非零なので$\det M\ne0$は
$\det\bigl((P^{\langle-1\rangle})^{\mathsf t}\bigr)=\det Q\ne0$
と同値である．これで(ii)$\Leftrightarrow$(iii)を得る．

最後に(iii)$\Rightarrow$(i)を証明する．
$S=\operatorname{Bl}_{p_1,\ldots,p_6}\mathbb P^2$は滑らかな可換三次曲面で，
$H=-K_S$はampleである．AOU\cite[\S\S4,8，特に(4.1)，(4.7)，(4.9)，(8.4)--(8.7)]{AOU}
の九線束$E_v$とそのrolled-up helix代数を取る．
その幾何学的potentialの二つの関係系列は，基底を適切に規格化すると
\begin{equation}\label{eq:geometric-relations}
 w^{\mathrm{geo}}_{1bc}=P_{cb},\qquad
 w^{\mathrm{geo}}_{2bc}=1\qquad(0\le b,c\le2)
\end{equation}
となる．これは次の初等的な式からも直接確認できる．
$u_b=(P_{0b}P_{1b},P_{1b}P_{2b},P_{2b}P_{0b})$，
$L_{b,c}(X)=\det[p_{b+1},p_{c+4},X]$とし，
$D_b(X)=(u_{b+1}\times u_{b+2})\cdot
(X_1X_2,X_2X_3,X_3X_1)$（添字を$3$で読む）とする．
$D_b$は$p_{b+1}$以外の五点を通るconicである．
また$M_b(X)=\det[p_{b+2},p_{b+3},X]$と書くと，
行列式の展開により全ての$c$について
\begin{equation}\label{eq:six-point-identities}
 \sum_{b=0}^2P_{cb}D_b(X)L_{b,c}(X)=0,\qquad
 \sum_{b=0}^2M_b(X)L_{b,c}(X)=0
\end{equation}
が成立する．各関係空間はAOU(4.2)により一次元であるため，
\eqref{eq:six-point-identities}が\eqref{eq:geometric-relations}を与える．
これらの多項式恒等式は独立の整数係数による厳密検算コードにも収めた．

\eqref{eq:geometric-relations}より，その二系列の比は
$w^{\mathrm{geo}}_{2bc}/w^{\mathrm{geo}}_{1bc}
=P_{cb}^{-1}=q_{bc}$である．
元の二系列と幾何学的二系列は各矢$y_{bc}$の非零rescalingで一致する．
残りの$w_{0bc}$は頂点$A_0$を除いた商$C$に影響しないから，
\[
 C(w)\simeq C\bigl(w^{\mathrm{geo}}\bigr)
\]
が道長次数付き代数として成り立つ．

幾何学的代数については，AOU(4.9)により
\begin{equation}\label{eq:geometric-rolled-orbit}
 B\bigl(w^{\mathrm{geo}}\bigr)
 \simeq\bigoplus_{v,v'}\bigoplus_{n\ge0}
 \operatorname{Hom}_S(E_v,E_{v'}(nH))
\end{equation}
がcut次数付き代数として成立し，$e=e_{0,0}$に対応する線束は
$E_{0,0}=\mathcal O_S$である．
任意の$v$について$m\gg0$を固定すれば
$E_v^\vee(mH)$は大域生成される．評価写像の核を$K_v$として
\[
 0\longrightarrow K_v\longrightarrow
 H^0(S,E_v^\vee(mH))\otimes\mathcal O_S
 \longrightarrow E_v^\vee(mH)\longrightarrow0
\]
を取る．これを$E_{v'}((n-m)H)$でtensorし，
Serre消滅$H^1(S,K_v\otimes E_{v'}((n-m)H))=0$（$n\gg0$）を使うと，
乗法写像
\begin{equation}\label{eq:serre-factorization}
 H^0(S,E_v^\vee(mH))\otimes
 H^0(S,E_{v'}((n-m)H))
 \twoheadrightarrow
 \operatorname{Hom}_S(E_v,E_{v'}(nH))
\end{equation}
が全射となる．左辺の積は必ず
$E_v\to\mathcal O_S(mH)\to E_{v'}(nH)$という，
構造頂点$e$を通る経路に対応する．頂点は有限個なので全ての$v,v'$に共通の
十分大きい$N$を取れば，
$B_n\subset BeB$（$n\ge N$）が従う．
cut次数零は有限次元で，各斉次成分も有限次元であるから，
$\dim_kC\bigl(w^{\mathrm{geo}}\bigr)<\infty$となる．
これで(i)が証明された．

最後に，標識付き六点のmoduliで座標三点$p_4,p_5,p_6$を固定すると，
残る自由度は$P$の各列のスカラーと座標三点を保つ対角$\mathrm{PGL}_3$で割ることである．
この商はちょうど$Q$の行・列rescaling商であり，
$P=(Q^{\langle-1\rangle})^{\mathsf t}$およびその逆は
全成分非零のトーラス上で正則だから，主張した開集合の同型が得られる．
\end{proof}

\begin{corollary}[AS heart と anticanonical heart の一致領域]
\label{cor:complete-as-ac}
十八係数非零の係数チャートで，標準AOU型のfolding，
反Serre作用のcut shiftとの同定，必要なnoetherian性を満たすAS対象に限る．
このとき自然なcorner restrictionが
\[
 \mathcal A_{\mathrm{AS}}\xrightarrow{\sim}
 \mathcal A_{\mathrm{ac}}=\qgr R(\mathcal O,\tau)
\]
という構造対象・反Serre作用と両立するabel同値であることと，
定理\ref{thm:general-position-iff}の十一条件は同値である．
\end{corollary}
\begin{proof}
定理\ref{thm:general-position-iff}と命題\ref{prop:corner}，
式\eqref{eq:orbit}を適用する．
\end{proof}

\begin{remark}
定理\ref{thm:general-position-iff}は，\emph{元の27係数potential全体が}
可換三次曲面から生じると主張するものではない．
$w_{0bc}$は任意でよく，幾何学的潜在係数と一致させる必要はない．
同定しているのは$A_0$除去後のcorner商である．
また，十一条件に違反した場合に任意の別のabel同値・導来同値が存在しないという主張でもない．
Rains heart，VdB heartとの比較は依然として独立の課題である．
\end{remark}

\section{一般位置における一様な消滅次数は七である}
\label{sec:uniform-c7}
本節は定理\ref{thm:general-position-iff}の有限次元性を定量化する．
次数とは特に指定しない限り道長次数を指し，cut次数と区別する．
有限体における既存の計算は証明の入力にしない．

\begin{theorem}[最良の一様消滅次数]\label{thm:uniform-c7}
標数零の代数閉体上，十八の係数$w_{1bc},w_{2bc}$が非零であり，
$Q=(w_{2bc}/w_{1bc})$とする．次は同値である：
\begin{enumerate}
\item $Q$の九つの$2\times2$小行列式，$\det Q$，
$\det Q^{\langle-1\rangle}$がすべて非零．
\item $C(w)=B(w)/B(w)e_{A_0}B(w)$が有限次元．
\item 道長次数七の成分$C(w)_7$が零．
\end{enumerate}
これらの条件の下で$C(w)_n=0$ $(n\ge7)$，
一方$\dim_kC(w)_6\ge8$である．従って四交比トーラス上
\[
 U=U_7=\{\Delta_{11}\ne0\},
\]
かつ七は最良の一様消滅次数である．$\Delta_{11}$は
十一の禁止divisorの定義多項式の積を表す．
\end{theorem}
\begin{proof}
(iii)$\Rightarrow$(ii)は$C$が矢と頂点で生成されることから従い，
(ii)$\Rightarrow$(i)は既存の必要条件である．
定理\ref{thm:general-position-iff}で六点を復元し，
$S=\operatorname{Bl}_{p_1,\ldots,p_6}\mathbb P^2$と置く．
$h$を直線の引き戻し，$e_i$を六つの例外曲線とし，
\[
\begin{aligned}
H&=3h-\sum_{i=1}^6e_i=-K_S, &
D_1&=-2h+\sum_{i=1}^6e_i,\\
D_2&=-h+e_1+e_2+e_3, &
q_j&=2h-\sum_{r\ne j}e_r \quad(4\le j\le6).
\end{aligned}
\]
$q_j$は五点を通るconicの真の変換で，
$q_j^2=-1$，$H\cdot q_j=1$．三blockの線束は
$\mathcal O(D_a)$ $(a=1,2)$，$\mathcal O(e_i)$ $(1\le i\le3)$，
$\mathcal O(h-e_j)$ $(4\le j\le6)$と，
構造頂点$A_0$の$\mathcal O_S$である．
AOU\cite[(4.9)]{AOU}の幾何的rolled-up helix代数を使う．
道長七の始点・終点は三種類（二十一組）だけであり，
対応する切断は
\[
\begin{array}{c|c|c}
\text{始点と終点}&\text{組数}&\text{対応するdivisor}\\\hline
 A_a\to B_i\ (a=1,2;\ 1\le i\le3)&6&2H+e_i-D_a\\
 B_i\to C_j\ (1\le i\le3;\ 4\le j\le6)&9&2H+h-e_i-e_j\\
 C_j\to A_a\ (4\le j\le6;\ a=1,2)&6&3H+D_a-h+e_j .
\end{array}
\]
それぞれの全切断が途中で$\mathcal O_S(mH)$を通ることを示す．

まず，四つの線束
\[
\begin{aligned}
L^-_1&=H-D_1=5h-2\sum e_r,&
L^-_2&=H-D_2=4h-2(e_1+e_2+e_3)-(e_4+e_5+e_6),\\
L^+_1&=H+D_1=h,&
L^+_2&=H+D_2=2h-e_4-e_5-e_6
\end{aligned}
\]
は，いずれも六本の互いに素な$(-1)$曲線を縮約する
$S\to\mathbb P^2$の直線束である．
実際，$L^-_1$は$q_1,\ldots,q_6$，
$L^-_2$は$h-e_1-e_2,h-e_1-e_3,h-e_2-e_3,q_4,q_5,q_6$，
$L^+_1$は$e_1,\ldots,e_6$，
$L^+_2$は$e_1,e_2,e_3,h-e_4-e_5,h-e_4-e_6,h-e_5-e_6$
を縮約する．従って$h^0(L_a^\pm)=3$．
$L^-_a\cdot e_i=2$ $(i\le3)$，
$L^+_a\cdot q_j=2$ $(j\ge4)$であり，
これらの$(-1)$曲線は縮約されず，平面の滑らかなconicへ
同型に写る．したがって
\begin{equation}\label{eq:c7-conic-restriction}
H^0(L^-_a)\xrightarrow{\sim}H^0(e_i,\mathcal O_{e_i}(2)),
\qquad
H^0(L^+_a)\xrightarrow{\sim}H^0(q_j,\mathcal O_{q_j}(2))
\end{equation}
は同型である．

\paragraph{$A_a\to B_i$（六組）．}
$N=2H+e_i-D_a$について
\[
0\to\mathcal O(2H-D_a)\to\mathcal O(N)
\to\mathcal O_{e_i}(2)\to0 .
\]
$3H-D_a$はampleなので
$H^1(2H-D_a)=0$（Kodaira消滅）．
左端の切断は$A_a\to\mathcal O(2H)\to B_i(2H)$を経由する．
$H^1(H)=0$から$H^0(H+e_i)\to H^0(e_i,\mathcal O_{e_i})$
は全射である．式\eqref{eq:c7-conic-restriction}を併用すると，
$A_a\to\mathcal O(H)\to B_i(2H)$を経由する切断は
$H^0(e_i,\mathcal O_{e_i}(2))$全体に制限される．
従って二種類の積が$H^0(N)$を張る．

\paragraph{$B_i\to C_j$（九組）．}
$L=h-e_j$はbase-point-free pencilである．
$3H-e_i-h+e_j$はnefかつbigである．
実際，二十七本の$(-1)$曲線との交点はすべて非負で，
自己交点は八である（del Pezzoの錐定理を用いる）．
Kawamata--Viehweg消滅で$H^1(2H-e_i-h+e_j)=0$．
pencilの評価列を$\mathcal O(2H-e_i)$でtensorして
\[
H^0(2H-e_i)\otimes H^0(h-e_j)
\twoheadrightarrow H^0(2H+h-e_i-e_j)
\]
を得る．右辺は全て$B_i\to\mathcal O(2H)\to C_j(2H)$を経由する．

\paragraph{$C_j\to A_a$（六組）．}
$N=3H+D_a-h+e_j=2H+D_a+q_j$とする．
$3H+D_a$はampleだから$H^1(2H+D_a)=0$で，
\[
0\to\mathcal O(2H+D_a)\to\mathcal O(N)
\to\mathcal O_{q_j}(2)\to0
\]
は切断について右端まで全射である．左端は
$q_j$の定義切断を掛けた
$C_j\to\mathcal O(H)\to A_a(3H)$の積で張られる．
$H^1(H)=0$と$(H+q_j)\cdot q_j=0$から
$H^0(H+q_j)\to H^0(q_j,\mathcal O_{q_j})$が全射である．
式\eqref{eq:c7-conic-restriction}より
$H^0(H+D_a)\to H^0(q_j,\mathcal O_{q_j}(2))$は同型．
従って$C_j\to\mathcal O(2H)\to A_a(3H)$の積は
右端の全切断を張り，先の積と合わせて$H^0(N)$を生成する．

以上より，全二十一組で道長七の全元が$BeB$に属し，
$C_7=0$である．

\paragraph{道長六が残ること．}
八頂点$A_1,A_2,B_1,B_2,B_3,C_4,C_5,C_6$の
各対角の道長六の全空間は
$H^0(2H)$で，Riemann--RochとKodaira消滅により次元十．
$A_a$では$H\cdot D_a=0$と$H$のamplenessから
$H^0(\pm D_a)=0$である．
従って構造頂点を経由する積は$\mathcal O(H)$経由に限られ，
\[
H^0(H-D_a)\otimes H^0(H+D_a)\longrightarrow H^0(2H)
\]
の像の次元は高々$3\cdot3=9$である．
$B_i$では$B_i\to\mathcal O$はなく，
$\mathcal O(H)$経由の積を$e_i$へ制限すると
$H^0(H-e_i)$の高々二次元の像となる．
$\mathcal O(2H)$経由の積は$e_i$上で零．
一方，$H^1(2H-e_i)=0$から$H^0(2H)$は
$H^0(e_i,\mathcal O_{e_i}(2))$（三次元）へ全射．
$C_j$では$C_j\to\mathcal O$はなく，
$\mathcal O(H)$経由の積は$q_j$上で零，
$\mathcal O(2H)$経由の積を$q_j$へ制限すると
$H^0(h-e_j)$の高々二次元の像となる．
$H^1(2H-q_j)=H^1(H+h-e_j)=0$なので
$H^0(2H)$は$H^0(q_j,\mathcal O_{q_j}(2))$（三次元）へ全射．
従って八頂点の対角商は各々非零で，
$\dim C_6\ge8$．
\end{proof}

\begin{corollary}[道長七の関係行列の根基]\label{cor:c7-radical}
$\Lambda=k[a^{\pm1},b^{\pm1},c^{\pm1},d^{\pm1}]$とし，
規格化したpotentialの道長七のJacobi関係行列を$M_7$，
その列数（六千八百四）の最大小行列式のイデアルを$I_7$とする．
このとき
\[
\sqrt{I_7}=\sqrt{(\Delta_{11})}\qquad\text{in }\Lambda .
\]
\end{corollary}
\begin{proof}
定理\ref{thm:uniform-c7}より，
$M_7$の列full rank locusは$\{\Delta_{11}\ne0\}$に等しい．
従って$V(I_7)=V(\Delta_{11})$であり，
Laurent多項式環上のHilbert Nullstellensatzによって根基が一致する．
これは根基の等式であり，非被約scheme構造や個々のminorの
因数分解を述べていない．
\end{proof}

\section{一般位置六点における corner の Hilbert 関数の一様性}
\label{sec:hilbert-constant}
\begin{theorem}[corner の完全な Hilbert 関数]\label{thm:hilbert-constant}
標数零の代数閉体上，十八の係数 $w_{1bc},w_{2bc}$ が非零で，
定理\ref{thm:general-position-iff}の十一の非零条件を満たすとする．
$C=B/BeB$ に道長 grading を入れると，六点配置には依存せず
\[
 \boxed{H_C(t)=8+21t+33t^2+38t^3+33t^4+21t^5+8t^6,\qquad
       \dim_k C=162.}
\]
特に定理\ref{thm:uniform-c7}の $\dim C_6\ge8$ は等号であり，
$C_0,\ldots,C_6$ の次元跳びは一般位置 locus に存在しない．
この結論は $B$ 全体が AS 正則であることを仮定しない．
\end{theorem}

\begin{proof}
まず定理\ref{thm:general-position-iff}により，
$C$ を一般位置の六点を吹き上げた滑らかな可換三次曲面 $S$ の
rolled-up helix 代数の同じ idempotent 商として計算できる．
既存の記号を保ち，$H=-K_S$，
\[
\begin{gathered}
 D_1=-2h+\sum_{r=1}^6e_r,\qquad D_2=-h+e_1+e_2+e_3,\\
 A_0=\mathcal O_S,\quad A_a=\mathcal O_S(D_a)\ (a=1,2),\quad
 B_i=\mathcal O_S(e_i)\ (1\le i\le3),\\
 C_j=\mathcal O_S(G_j),\qquad G_j=h-e_j\ (4\le j\le6),\qquad
 q_j=H-G_j=2h-\sum_{r\ne j}e_r.
\end{gathered}
\]
ここで $q_j$ は五点を通る conic の真の変換である $(-1)$ 曲線の類を表す．
$A,B,C$ の三 block にレベル $0,1,2$ を与えると，
道長 $n$ の始点 $s$，終点 $t$ の全 Hom は，
AOU\cite[(4.1), (4.9)]{AOU}の同定により
\begin{equation}\label{eq:hilbert-block}
 B_n(s,t)=H^0\!\left(S,\mathcal O_S\left(
 \left\lfloor\frac{\ell(s)+n}{3}\right\rfloor H+E_t-E_s\right)\right)
 \quad(\ell(t)\equiv\ell(s)+n\pmod3),
\end{equation}
ただし $E_{A_a}=D_a$，$E_{B_i}=e_i$，$E_{C_j}=G_j$ である．
AOU の helix の acyclicity と曲面の Riemann--Roch
$\chi(\mathcal O(D))=1+D(D+H)/2$ を使うと，
$e=A_0$ を除く八頂点の全 Hom の各次数の合計は次の表の中央左列になる．
\[
\begin{array}{c|rrrrrrrr}
\text{道長 }n&0&1&2&3&4&5&6&7\\\hline
\dim B_n\text{（八頂点間）}&8&21&42&74&105&147&206&252\\
\dim (BeB)_n\text{（同上）}&0&0&9&36&72&126&198&252\\
\dim C_n&8&21&33&38&33&21&8&0
\end{array}
\]
以下では中央右列，すなわち構造頂点を\emph{途中で経由する}
積の像の次元を各次数について証明する．
多重経由があるときは，すべての経由位置に対応する像の和を取る．

必要な切断数を列挙する（$a=1,2$，$i=1,2,3$，$j=4,5,6$）：
\begin{equation}\label{eq:hilbert-basic-h0}
\begin{gathered}
 h^0(H)=4,\quad h^0(2H)=10,\quad h^0(H\pm D_a)=3,\\
 h^0(H-e_i)=2,\quad h^0(H+e_i)=5,\quad h^0(2H-e_i)=7,\\
 h^0(G_j)=2,\quad h^0(q_j)=1,\quad
 h^0(H+G_j)=7,\quad h^0(H+q_j)=5.
\end{gathered}
\end{equation}
$H-e_i$ と $G_j$ は base-point-free pencil である．
$H\pm D_a$ は定理\ref{thm:uniform-c7}で使った平面への blow-down 線束である．

\paragraph{道長二と三．}
道長二では $C_j\to B_i$ の九組のみが $e$ を経由し，
その像は $H^0(q_j)H^0(e_i)$ の一次元である．合計九次元．
道長三では $A_a\to A_b$ に経由射はなく，
$B_i\to B_k$ と $C_j\to C_k$ の九組ずつで，
それぞれ $H^0(H-e_i)H^0(e_k)$ および
$H^0(q_j)H^0(G_k)$ が二次元の像を与える．合計三十六次元．

\paragraph{道長四．}
全二十一組の全 Hom はそれぞれ五次元である．
$A_a\to B_i$（六組）では $H^0(H-D_a)H^0(e_i)$ の三次元が経由する．
$B_i\to C_j$（九組）では
$H^0(H-e_i)H^0(G_j)$ が経由する．後者の pencil trick の核は
$H^0(q_j-e_i)$ であり，これは零である．実際 $q_j\cdot e_i=1$ と
$q_j^2=-1$ から，そのような effective divisor は $q_j$ を成分として含むが，
差は $-e_i$ となって effective でない．従って積の像は四次元．
$C_j\to A_a$（六組）では $H^0(q_j)H^0(H+D_a)$ の三次元が経由する．
この三種類を加えると $6\cdot3+9\cdot4+6\cdot3=72$ である．

\paragraph{道長五．}
全二十一組の全 Hom はそれぞれ七次元である．
$A_a\to C_j$（六組）では
$H^0(H-D_a)H^0(G_j)$ の像を取る．$G_j$ の pencil trick の核は
$H^0(H-D_a-G_j)$ であり，零である．
$a=1$ なら $H-D_1-G_j=2q_j-e_j$ で，
$q_j$ との交点数は $-2$ だから $q_j$ が二重の固定成分となり，
残りが $-e_j$ となって矛盾する．$a=2$ なら
$H-D_2-G_j=3h-2(e_1+e_2+e_3)-\sum_{r=4,r\ne j}^6e_r$．
これは互いに素な三本の $(-1)$ 曲線
$h-e_1-e_2$，$h-e_1-e_3$，$h-e_2-e_3$ とそれぞれ交点数 $-1$ であり，
三本を引くと $-\sum_{r=4,r\ne j}^6 e_r$ となり effective でない．
よって積の像は $3\cdot2=6$ 次元．
$B_i\to A_a$（六組）では
$H^0(H-e_i)H^0(H+D_a)$ の像を取る．$H-e_i$ の pencil trick の核は
$H^0(D_a+e_i)=0$（$h$ との交点が負）なので六次元．
$C_j\to B_i$（九組）は $e$ を二通りの位置で経由し，
$H^0(H+q_j+e_i)$ の中で
\[
 q_j\,H^0(H+e_i)+e_i\,H^0(H+q_j)
\]
を得る．ここで曲線の記号はその定義切断を掛ける作用も表す．
二つの部分空間は共に五次元であり，$q_j$ と $e_i$ に共通成分がないので
交叉は $q_je_i\,H^0(H)$ の四次元である．よって和は六次元．
従って経由射の合計は $(6+6+9)\cdot6=126$ 次元である．

\paragraph{道長六：$A_a\to A_b$（四組）．}
この場合の経由射は一種類で，次の積の像である：
\begin{equation}\label{eq:two-blowdowns-injection}
 H^0(H-D_a)\otimes H^0(H+D_b)
 \longrightarrow H^0(2H+D_b-D_a).
\end{equation}
この積写像が全ての $a,b=1,2$ に対して単射であることを示す．
$L=H-D_a$，$M=H+D_b$ とする．いずれも $S\to\mathbb P^2$ の
blow-down を定め，切断空間は三次元である．
$L$ が縮約する六本の $(-1)$ 曲線のうち，$M\cdot E=2$ となる曲線は
$a=b$ で六本，$a\ne b$ で三本（後者は $q_4,q_5,q_6$）ある．
$L$ により縮約されるそれら三本の像は $\mathbb P^2$ で非共線である．
もし共線なら，その直線の真の変換は $H$ と交点零の有効曲線となり，
$H$ の ampleness に矛盾する．
一方，$M\cdot E=2$ なら，$M$ による blow-down は $E$ を
滑らかな平面 conic に同型に写す．したがって
\[
 H^0(M)\xrightarrow{\sim}H^0(E,M|_E)\simeq H^0(\mathbb P^1,\mathcal O(2))
\]
である．
仮に \eqref{eq:two-blowdowns-injection} の核に $\sum_{u=1}^3s_u\otimes t_u$
（$s_u$ は $H^0(L)$ の基底）があれば，三本の曲線への制限から
$\sum_{u=1}^3s_u(p_E)t_u=0$ が $H^0(M)$ で成り立つ．
三つの $p_E$ が非共線なので，三つの評価汎関数は独立であり，
全ての $t_u$ が零である．これで積写像の単射性が従う．
従って全ての四組で経由射の像は九次元である．

\paragraph{道長六：$B_i\to B_k$（九組）．}
$H^0(H-e_i)H^0(H+e_k)$ の像は，
$H-e_i$ の pencil trick と $h^0(e_i+e_k)=1$ より九次元である．
$i\ne k$ では全 Hom が九次元なのでこれで終わる．
$i=k$ のときは二番目の経由位置から $e_i\,H^0(2H-e_i)$ も得られるが，
$H^0(H-e_i)H^0(H)\twoheadrightarrow H^0(2H-e_i)$ が成立する．
実際，pencil trick の核は $H^0(e_i)$ の一次元で，
$2\cdot4-1=7=h^0(2H-e_i)$ だからである．
従って二番目の積は最初の像に含まれ，九次元のままである．

\paragraph{道長六：$C_j\to C_k$（九組）．}
$H^0(H+q_j)H^0(G_k)$ の像は，$G_k$ の pencil trick と
$h^0(q_j+q_k)=1$ より九次元である．
$j\ne k$ では全 Hom が九次元なのでこれで終わる．
$j=k$ の場合，もう一つの経由位置から $q_j\,H^0(H+G_j)$ も得るが，
$H^0(H)H^0(G_j)\twoheadrightarrow H^0(H+G_j)$ が成立する．
実際，pencil trick の核は $H^0(q_j)$ の一次元で，
$4\cdot2-1=7=h^0(H+G_j)$ である．
従って後者の積も先の九次元の像に含まれる．

以上より道長六の経由射の次元は
$(4+9+9)\cdot9=198$ である．
表中の全 Hom の次元 $206$ との差は八であり，$\dim C_6=8$ が従う．
道長七以降の消滅は定理\ref{thm:uniform-c7}による．
表を合計して Hilbert 関数と $\dim C=162$ を得る．
\end{proof}

\begin{corollary}[162次元 algebra の有限平坦族]
\label{cor:flat-corner-162}
標数零の代数閉体 $k$ 上の四交比トーラスにおける十一条件の補開集合
$U=\{\Delta_{11}\ne0\}$ 上，規格化した universal quiver potential の
$e_{A_0}$-corner 商を $\mathcal C_U$ とする．このとき
$\mathcal C_U$ は階数 $162$ の有限局所自由な
$\mathcal O_U$-代数で，その道長斉次部分は順に
\[
 (8,21,33,38,33,21,8)
\]
の階数を持ち，道長七以上は零である．任意の $U$ 上の基底変更と可換する．
\end{corollary}
\begin{proof}
各道長 $n$ の成分は有限個の道を基底とする自明ベクトル束を，
全ての $u(\partial W)v$ からなる有限個の関係によって割った
有限表示 $\mathcal O_U$ 加群である．
定理\ref{thm:hilbert-constant}より $0\le n\le6$ では
全幾何学的ファイバーの次元が一定であるから，
対応する関係行列は一定階数を持ち，商は局所自由で基底変更と可換する．
$n=7$ では全ファイバーが零なので，Nakayamaの補題で
$\mathcal C_{U,7}=0$．頂点と矢による生成性から高次数も零になる．
有限個の斉次部分を足し合わせて主張を得る．
\end{proof}

\begin{remark}[結果の対象範囲]\label{rem:hilbert-status}
証明は四交比の十一条件による一般位置六点の再構成と，
可換三次曲面の幾何学的 rolled-up helix の Hom 同定に基づく．
$w_{0bc}$ は $C$ に影響しないので任意でよい．
しかし $27$ 係数の任意の potential が AS-regular/noetherian であること，
Rains heart や VdB heart と同一であること，導来圏の全域での同値は
ここからは結論しない．
\end{remark}

\paragraph{根壁との関係．}十一判別因子上の古典的収縮と弱楕円gluingの核については主論文\cite{Hearts}を参照する．
\section{非空性の明示的な証明書}
\subsection{六点からpotentialを復元する}
次の点を$\mathbb P^2$に取る：
\[
 p_1=(1:0:0),\ p_2=(0:1:0),\ p_3=(0:0:1),\quad
 p_4=(1:1:1),\ p_5=(1:2:3),\ p_6=(1:4:2).
\]
全三点行列式は非零であり，二次単項式
$(x^2,xy,xz,y^2,yz,z^2)$の六点評価行列式は$-12$である．
従って三点は共線でなく，六点はconic上にない．そのblow-up $S$は滑らかなcubicである．

AOUのfoundationの三blockを
\[
 \begin{gathered}
 (\OO,\OO(-2h+\textstyle\sum e_i),\OO(-h+e_1+e_2+e_3)),\\
 (\OO(e_1),\OO(e_2),\OO(e_3)),\qquad
 (\OO(h-e_4),\OO(h-e_5),\OO(h-e_6))
 \end{gathered}
\]
とする\cite[(4.1)]{AOU}．
$\ell_{ij}$を$p_i,p_j$を通る直線，$q_i$を$p_i$以外の五点を通るconicとする．
係数は原始整数列で最初の非零成分を正として規格化する．
$A_a\to B_b\to C_c$の三つの合成切断は，$i=b+1,j=c+4$と置くと
\[
 \ell_{ij},\qquad q_i\ell_{ij},\qquad
 \ell_{\{1,2,3\}\setminus\{i\}}\ell_{ij}
\]
で表される（左から$a=0,1,2$）．
各$a,c$に対する三切断の唯一の関係を同じ原始規格化で取り，$w_{abc}$とする．
これはAOUのrelation-extension potentialの直接構成\cite[(4.7)]{AOU}である．
得られる行列$W_a=(w_{abc})_{b,c}$は
\begin{equation}\label{eq:witness}
 W_0=\begin{pmatrix}1&1&1\\-1&-2&-2\\1&3&1\end{pmatrix},\quad
 W_1=\begin{pmatrix}2&1&1\\-1&2&-4\\-1&-3&1\end{pmatrix},\quad
 W_2=\begin{pmatrix}1&1&2\\-1&-1&-1\\1&1&1\end{pmatrix}.
\end{equation}
各関係は多項式係数として厳密に検証する．Euler行列だけからの推測ではない．
この可換点でのAS性はAOUのhelixの既知結果を用い，有限rankから推論しない．
noetherian性は，rolled algebraが反標準section環上の有限加群であることから従う．
実際，九線束の各Hom束の十分高次のsectionは有限生成であり，低次数も有限次元である．

\subsection{商の全関係行列}
$A_0$を除くと，頂点8，矢21，二次関係21が残る．
残る全関係を全てのprefixとsuffixで掛け，始点・終点別に行列を作る．
\begin{proposition}[厳密な有限計算]\label{prop:certificate}
式\eqref{eq:witness}のpotentialに対し，道長gradingで
\[
 H_{B/BeB}(t)=8+21t+33t^2+38t^3+33t^4+21t^5+8t^6.
\]
特に$w_*\in U_7$であり，$U_7$は非空である．
\end{proposition}
\begin{proof}
再現コード\texttt{computations/verify\_cubic\_corner.py}は，点から式\eqref{eq:witness}を復元し，
全Jacobi関係を微分から構成する．$C_n$のpresentation行列は全ての
$u(\partial W)v$を列挙しており，二次生成イデアルの次数$n$部分を尽くす．
有理数上の分数演算による行消去で次を得る：
\begin{center}
\begin{tabular}{c|rrrrrrrr}
道長$n$&0&1&2&3&4&5&6&7\\\hline
道の数&8&21&54&144&378&972&2592&6804\\
$\dim C_n$&8&21&33&38&33&21&8&0
\end{tabular}
\end{center}
同じ行列を素数$1000003$で還元した別の消去でも同じrankを得る．
7次は21個の始終点blockの全てで列full rankであり，計6804列を消去する．
finite-fieldでのfull rankだけでも整数minorの非零を証明するが，低次数の正確な次元には
有理数計算を用いている．浮動小数点は用いない．
$C$は道長1で生成されるので$C_7=0$から全高次数が零となり，表示が全Hilbert級数となる．
総次元は$162$である．証明書の行列生成と検算は計算機に依存するが，
新しいparameterについてのAS性をこの検算で置き換えてはいない．
\end{proof}

以上により，標識と構造頂点を固定したnoetherian AS領域で，AS heartとanticanonical heartが
一致する明示的な非空開集合を得た．係数は全て非零なのでpotentialの係数torusとも交わる．
矢の非零rescalingは$e$を保つため，有限次元性と$U_N$はその作用で不変である．
従って標識を保つ商にも同じ開条件が降りる．

\section{HRS比較を進める符号判定}
主論文\cite{Hearts}の振幅1比較を$F:D^b(\cA)\simeq D^b(\cB)$と書き，
$F(\cA)\subset D^{[0,1]}(\cB)$，$F^{-1}(\cB)\subset D^{[-1,0]}(\cA)$とする．
$(T,U)=(\ker H^1F,\ker H^0F)$はtorsion pairである．
その理由を記す．$M\in\cA$について$B_i=H^iFM$とし，truncation triangleを
$G=F^{-1}$で移すと$GB_0\to M\to GB_1[-1]$となる．
振幅条件から$H^{-1}GB_0=H^0GB_1=0$であり，
$0\to GB_0\to M\to H^{-1}GB_1\to0$を得る．
端の二項はそれぞれ$T,U$に属し，$\Hom(T,U)=0$は
$F(T)\subset\cB$，$F(U)\subset\cB[-1]$から従う．
従って$G(\cB)=\langle U[1],T\rangle$である．
\begin{proposition}\label{prop:sign}
$F(e)=e'$，$F\tau=\tau'F$で，両作用が対応するheartを保つとする．
$(e',\tau')$がAZ-ampleで高次Serre消滅を満たすと仮定する．
$0\ne M\in\cA$が$\tau M\simeq M$かつ単純なら，$\chi(e,M)\ne0$であり，
\[
 M\in T\ \Longleftrightarrow\ \chi(e,M)>0,\qquad
 M\in U\ \Longleftrightarrow\ \chi(e,M)<0.
\]
$M\in U$なら$\Ext^i(e,M)$は$i=1$だけで非零となり，
$\dim\Ext^1(e,M)=-\chi(e,M)$である．
\end{proposition}
\begin{proof}
単純性とtorsion分解により$M$は$T$または$U$に属する．
$N=H^0FM$または$N=H^1FM$は非零で$\tau'$固定である．
高次消滅と固定性から$\Ext^{>0}(e',N)=0$．
もし$\Hom(e',N)=0$なら全twistからのHomが零となりAZ生成性に反する．
従って$h=\dim\Hom(e',N)>0$．$FM=N$ならEuler値は$h$，
$FM=N[-1]$なら$-h$であり，後者のRHomは次数1に集中する．
\end{proof}
この符号判定は，VdBの各段階でpointednessと反標準作用の適合性を検証した後に適用する．
共鳴式だけから必要な単純対象の存在を作ったことにはならない．
一般の非単純対象について，Euler値が負というだけで対象全体が$U$に入るとはしない．

\section{結果の範囲と残る問題}
定理\ref{thm:open}は一般のcorner原理の証明と，cubicの有限表示での明示化である．
命題\ref{prop:certificate}の証明は有理数・有限体の厳密計算を含む．
外部独立査読と新規性の全面調査は未実施．

反標準環はAS側の指定対象$O$から作ったものなので，
AOU \S9のmutation後の先頭対象がVdB模型の$\pi F$と同一であるか，
および既知の特定の導来同値がこのcorner同値に一致するかは別問題である．
今回そこを仮定なしに解決したとはしない．
また$U_7$の外が全て非一致であるわけではなく，より高次の$U_N$に入ることがある．
Rains heartとVdB heart（後者はliteral blow-upに由来する），共鳴ごとのtorsion pair，特異境界は引き続き比較する．
数値的共鳴の稠密性から，これらの導来非同値の密度を推論してはいけない．



\begin{thebibliography}{99}

\bibitem{RainsBirational}
E.~M.~Rains, \emph{The birational geometry of noncommutative surfaces},
arXiv:1907.11301v1, Theorem 4.30, \S7, Theorem 9.10,
Corollaries 10.12, 10.16, Proposition 10.13, \S11.
\url{https://arxiv.org/abs/1907.11301}.
\bibitem{RainsFiltered}
E.~M.~Rains, \emph{Filtered deformations of elliptic algebras},
arXiv:2107.13540v2 (2022), Corollary 3.8, \S4, Proposition 5.1.
\url{https://arxiv.org/abs/2107.13540}.
\bibitem{LiuWangZigzag} J.~Liu and Z.~Wang,
\emph{$A_\infty$-deformations of zigzag algebras via Ginzburg dg algebras},
J. Algebra 677 (2025), 360--371, Theorem 1.1; arXiv:2301.06757v3.
\url{https://arxiv.org/abs/2301.06757}.
\bibitem{AZ} M.~Artin, J.~J.~Zhang, \emph{Noncommutative projective schemes},
Adv. Math. 109 (1994), 228--287.
\bibitem{RSSResonance} D.~Rogalski, S.~J.~Sierra, J.~T.~Stafford,
\emph{Noncommutative Blowups of Elliptic Algebras},
arXiv:1308.2216v2 (2015), Proposition 2.4.
\url{https://arxiv.org/abs/1308.2216}.
\bibitem{YZResonance} A.~Yekutieli, J.~J.~Zhang,
\emph{Rings with Auslander Dualizing Complexes},
J. Algebra 213 (1999), 1--51, Theorems 4.5, 5.1(1), 5.13;
arXiv:math/9804005v3. \url{https://arxiv.org/abs/math/9804005}.
\bibitem{OrlovResonance} D.~Orlov,
\emph{Smooth and proper noncommutative schemes and gluing of DG categories},
arXiv:1402.7364v5 (2015), Theorems 3.18, 3.25, Corollary 3.19.
\url{https://arxiv.org/abs/1402.7364}.

\bibitem{AOU} T.~Abdelgadir, S.~Okawa, K.~Ueda,
\emph{Compact moduli of marked noncommutative cubic surfaces},
arXiv:2404.00175v1 (2024), \S\S3--4, 9.
\url{https://arxiv.org/abs/2404.00175}.
\bibitem{OU} S.~Okawa, K.~Ueda,
\emph{Marked Artin--Schelter surfaces of del Pezzo types},
arXiv:2508.04986v1 (2025), Definitions 1.1, 5.2 and \S6.
\url{https://arxiv.org/abs/2508.04986}.
\bibitem{VdB} M.~Van den Bergh,
\emph{Blowing up of non-commutative smooth surfaces},
Mem. Amer. Math. Soc. 154 (2001), no.~734;
arXiv:math/9809116v1, Theorems 9.1.9, 11.1.3.
\url{https://arxiv.org/abs/math/9809116}.
\bibitem{BlancLamy} J.~Blanc, S.~Lamy,
\emph{Weak Fano threefolds obtained by blowing-up a space curve and construction of Sarkisov links},
Proc. London Math. Soc. (3) 105 (2012), 1047--1075, Introduction and \S4.
\url{https://doi.org/10.1112/plms/pds023}.
\bibitem{Dolgachev} I.~V.~Dolgachev,
\emph{Classical Algebraic Geometry: A Modern View},
Cambridge University Press (2012), Theorems 8.3.2 and 8.3.4.
\bibitem{Hearts} 本研究，\emph{非可換del Pezzo曲面のheart比較}，
\texttt{papers/dp-hearts.tex}，2026年11月17日．
HRS表示と一般corner判定の現行証明．本稿は有限生成性を道長判定から導き，完全Hilbert関数を記述する．

\end{thebibliography}
\end{document}
